
% LaTeX Beamer file automatically generated from DocOnce
% https://github.com/doconce/doconce

%-------------------- begin beamer-specific preamble ----------------------

\documentclass{beamer}

\usetheme{red_plain}
\usecolortheme{default}

% turn off the almost invisible, yet disturbing, navigation symbols:
\setbeamertemplate{navigation symbols}{}

% Examples on customization:
%\usecolortheme[named=RawSienna]{structure}
%\usetheme[height=7mm]{Rochester}
%\setbeamerfont{frametitle}{family=\rmfamily,shape=\itshape}
%\setbeamertemplate{items}[ball]
%\setbeamertemplate{blocks}[rounded][shadow=true]
%\useoutertheme{infolines}
%
%\usefonttheme{}
%\useinntertheme{}
%
%\setbeameroption{show notes}
%\setbeameroption{show notes on second screen=right}

% fine for B/W printing:
%\usecolortheme{seahorse}

\usepackage{pgf}
\usepackage{graphicx}
\usepackage{epsfig}
\usepackage{relsize}

\usepackage{fancybox}  % make sure fancybox is loaded before fancyvrb

\usepackage{fancyvrb}
\usepackage{minted} % requires pygments and latex -shell-escape filename
%\usepackage{anslistings}
%\usepackage{listingsutf8}

\usepackage{amsmath,amssymb,bm}
%\usepackage[latin1]{inputenc}
\usepackage[T1]{fontenc}
\usepackage[utf8]{inputenc}
\usepackage{colortbl}
\usepackage[english]{babel}
\usepackage{tikz}
\usepackage{framed}
% Use some nice templates
\beamertemplatetransparentcovereddynamic

% --- begin table of contents based on sections ---
% Delete this, if you do not want the table of contents to pop up at
% the beginning of each section:
% (Only section headings can enter the table of contents in Beamer
% slides generated from DocOnce source, while subsections are used
% for the title in ordinary slides.)
\AtBeginSection[]
{
  \begin{frame}<beamer>[plain]
  \frametitle{}
  %\frametitle{Outline}
  \tableofcontents[currentsection]
  \end{frame}
}
% --- end table of contents based on sections ---

% If you wish to uncover everything in a step-wise fashion, uncomment
% the following command:

%\beamerdefaultoverlayspecification{<+->}

\newcommand{\shortinlinecomment}[3]{\note{\textbf{#1}: #2}}
\newcommand{\longinlinecomment}[3]{\shortinlinecomment{#1}{#2}{#3}}

\definecolor{linkcolor}{rgb}{0,0,0.4}
\hypersetup{
    colorlinks=true,
    linkcolor=linkcolor,
    urlcolor=linkcolor,
    pdfmenubar=true,
    pdftoolbar=true,
    bookmarksdepth=3
    }
\setlength{\parskip}{0pt}  % {1em}

\newenvironment{doconceexercise}{}{}
\newcounter{doconceexercisecounter}
\newenvironment{doconce:movie}{}{}
\newcounter{doconce:movie:counter}

\newcommand{\subex}[1]{\noindent\textbf{#1}}  % for subexercises: a), b), etc

\logo{{\tiny \copyright\ 1999-2023, Morten Hjorth-Jensen. Released under CC Attribution-NonCommercial 4.0 license}}

%-------------------- end beamer-specific preamble ----------------------

% Add user's preamble




% insert custom LaTeX commands...

\raggedbottom
\makeindex

%-------------------- end preamble ----------------------

\begin{document}

% matching end for #ifdef PREAMBLE

\newcommand{\exercisesection}[1]{\subsection*{#1}}



% ------------------- main content ----------------------



% ----------------- title -------------------------

\title{Week 34: Introduction to the course, Logistics and Practicalities}

% ----------------- author(s) -------------------------

\author{Morten Hjorth-Jensen\inst{1,2}}
\institute{Department of Physics and Center for Computing in Science Education, University of Oslo, Norway\inst{1}
\and
Department of Physics and Astronomy and Facility for Rare Isotope Beams, Michigan State University, USA\inst{2}}
% ----------------- end author(s) -------------------------

\date{Week 34, August 21-25, 2021
% <optional titlepage figure>
\ \\ 
{\tiny \copyright\ 1999-2023, Morten Hjorth-Jensen. Released under CC Attribution-NonCommercial 4.0 license}
}

\begin{frame}[plain,fragile]
\titlepage
\end{frame}

\begin{frame}[plain,fragile]
\frametitle{Overview of first week}

\begin{enumerate}
\item The sessions on Tuesdays and Wednesdays last four hours for each group (four groups in total) and will include lectures in a flipped mode (promoting active learning) and work on exercices and projects.

\item The sessions will begin with lectures, discussions, questions and answers about the material to be covered every week.

\item There are four groups:
\begin{itemize}

  \item Tuesdays 815am-12pm and 1215pm-4pm

  \item Wednesdays 815am-12pm and 1215pm-4pm.
\end{itemize}

\noindent
\end{enumerate}

\noindent
Please sign up as soon as
possible for one of the groups. Max capacity per group is 30-40
participants. The labs are also available till 6pm Tuesdays and Wednesdays. Videos and learning material with reading suggestions will be made available before each week starts.

\begin{enumerate}
\item On Thursdays we have a regular lecture. These lectures start at 1215pm and end at 2pm and serve the aims of giving an overview over various topics. 
\end{enumerate}

\noindent
The first week we start with simple linear regression, a repetition of
linear algebra and elements of statistics needed for the course.
\end{frame}

\begin{frame}[plain,fragile]
\frametitle{Schedule first week}

\begin{block}{}
\begin{itemize}
  \item August 22: Presentation of the course, aims and content. Introduction to software and repetition of Python Programming, linear algebra and basic elements of statistics. Please select group.

  \item August 23: Presentation of the course, aims and content. Introduction to software and repetition of Python Programming, linear algebra and basic elements of statistics. Please select group.

  \item August 24: Lecture: Linear regression, examples and theory 
\end{itemize}

\noindent
\end{block}
\end{frame}

\begin{frame}[plain,fragile]
\frametitle{Reading Recommendations}

\begin{block}{}
For the reading assignments we use the following abbreviations:
\begin{itemize}
\item GBC: \href{{https://www.deeplearningbook.org/}}{Goodfellow, Bengio, and Courville, Deep Learning}

\item CMB: Christopher M. Bishop, Pattern Recognition and Machine Learning

\item HTF: Hastie, Tibshirani, and Friedman, The Elements of Statistical Learning

\item AG: Aurelien Geron, Hands‑On Machine Learning with Scikit‑Learn and TensorFlow

\item KM: \href{{https://probml.github.io/pml-book/book1.html}}{Kevin Murphy, Probabilistic Machine Learning}
\end{itemize}

\noindent
Reading recommendations this week: Refresh linear algebra, GBC chapters 1 and 2. CMB sections 1.1 and 3.1. HTF chapters 2 and 3. Install scikit-learn. See lecture notes for week 34 at \href{{https://compphysics.github.io/MachineLearning/doc/web/course.html}}{\nolinkurl{https://compphysics.github.io/MachineLearning/doc/web/course.html}} (these notes).
\end{block}
\end{frame}

\begin{frame}[plain,fragile]
\frametitle{Lectures and ComputerLab}

\begin{block}{}
\begin{itemize}
  \item The sessions on Tuesdays and Wednesdays last four hours and will include partly lectures in a flipped mode (promoting active learning) and work on exercices and projects.

  \item Thursdays: regular lectures (12.15pm-2pm) 

  \item Weekly reading assignments and videos needed to solve projects and exercises.

  \item Weekly exercises when not working on projects. You can hand in exercises if you want.

  \item Detailed lecture notes, exercises, all programs presented, projects etc can be found at the homepage of the course.

  \item Weekly plans and all other information are on the official webpage.

  \item No final exam, three projects that are graded and have to be approved. 
\end{itemize}

\noindent
\end{block}
\end{frame}

\begin{frame}[plain,fragile]
\frametitle{Communication channels}

\begin{itemize}
\item Chat and communications via \href{{canvas.uio.no}}{\nolinkurl{canvas.uio.no}}

\item \textbf{Slack} channel: machinelearninguio.slack.com
\end{itemize}

\noindent
\end{frame}

\begin{frame}[plain,fragile]
\frametitle{Course Format}

\begin{block}{}
\begin{itemize}
  \item Three compulsory projects. Electronic reports only using \href{{https://www.uio.no/english/services/it/education/canvas/}}{Canvas} to hand in projects and \href{{https://git-scm.com/}}{git} as version control software and \href{{https://github.com/}}{GitHub} for repository (or \href{{https://about.gitlab.com/}}{GitLab}) of all your material.

  \item Evaluation and grading: The three projects are graded and each counts 1/3 of the final mark. No final written or oral exam.
\begin{enumerate}

   \item For the last project each group/participant submits a proposal or works with suggested (by us) proposals for the project.

   \item If possible, we would like to organize the last project as a workshop where each group makes a poster and presents this to all other participants of the course

   \item Poster session where all participants can study and discuss the other proposals.

   \item Based on feedback etc, each group finalizes the report and submits for grading. 

\end{enumerate}

\noindent
  \item Python is the default programming language, but feel free to use C/C++ and/or Fortran or other programming languages. All source codes discussed during the lectures can be found at the webpage and \href{{https://github.com/CompPhysics/MachineLearning/tree/master/doc/Programs}}{github address} of the course. 
\end{itemize}

\noindent
\end{block}
\end{frame}

\begin{frame}[plain,fragile]
\frametitle{Teachers}

\begin{block}{}

\textbf{Teachers :}
\begin{itemize}
\item Morten Hjorth-Jensen, morten.hjorth-jensen@fys.uio.no
\begin{itemize}

  \item \textbf{Phone}: +47-48257387

  \item \textbf{Office}: Department of Physics, University of Oslo, Eastern wing, room FØ470 

  \item \textbf{Office hours}: \emph{Anytime}!  Individual or group office hours can be arranged either in person or  via zoom. Feel free to send an email for planning. 

\end{itemize}

\noindent
\item Ida Torkjellsdatter Storehaug, i.t.storehaug@fys.uio.no

\item Fahimeh Najafi, fahim.n.20@gmail.com

\item Mia-Katrin Ose Kvalsund, m.k.o.kvalsund@fys.uio.no

\item Karl Henrik Fredly, k.h.fredly@fys.uio.no

\item Adam Jakobsen, adam.jakobsen@fys.uio.no

\item Daniel Haas Beccatini Lima, d.h.b.lima@fys.uio.no
\end{itemize}

\noindent
\end{block}
\end{frame}

\begin{frame}[plain,fragile]
\frametitle{Deadlines for projects (tentative)}

\begin{block}{}

\begin{enumerate}
\item Project 1: October 9 (available September 4) graded with feedback)

\item Project 2: November 6 (available October 6, graded with feedback)

\item Project 3: December 11  (available November 10, graded with feedback)
\end{enumerate}

\noindent
Projects are handed in using \textbf{Canvas}. We use Github as repository for codes, benchmark calculations etc. Comments and feedback on projects only via \textbf{Canvas}. 

\end{block}
\end{frame}

\begin{frame}[plain,fragile]
\frametitle{Recommended textbooks}

\begin{enumerate}
\item The lecture notes are collected as a jupyter-book at \href{{https://compphysics.github.io/MachineLearning/doc/LectureNotes/_build/html/intro.html}}{\nolinkurl{https://compphysics.github.io/MachineLearning/doc/LectureNotes/_build/html/intro.html}}.
\end{enumerate}

\noindent
In addition to the lecture notes, we recommend the books of Bishop and Goodfellow et al.~We will follow these texts closely and the weekly reading assignments refer to these two texts. The text by Hastie et al is also widely used in the Machine Learning community. Finally, we also recommend the hands-on text by Geron, see below.

\begin{enumerate}
\item Christopher M. Bishop, Pattern Recognition and Machine Learning, Springer, \href{{https://www.springer.com/gp/book/9780387310732}}{\nolinkurl{https://www.springer.com/gp/book/9780387310732}}. 

\item Ian Goodfellow, Yoshua Bengio, and Aaron Courville. The different chapters are available for free at \href{{https://www.deeplearningbook.org/}}{\nolinkurl{https://www.deeplearningbook.org/}}. Chapters 2-14 are highly recommended. The lectures follow to a larg extent this text. The weekly plans will include reading suggestions from these two textbooks.
\end{enumerate}

\noindent
Additional textbooks:

\begin{enumerate}
\item Trevor Hastie, Robert Tibshirani, Jerome H. Friedman, The Elements of Statistical Learning, Springer, \href{{https://www.springer.com/gp/book/9780387848570}}{\nolinkurl{https://www.springer.com/gp/book/9780387848570}}. This is a well-known text and serves as additional literature.

\item Aurelien Geron, Hands‑On Machine Learning with Scikit‑Learn and TensorFlow, O'Reilly, \href{{https://www.oreilly.com/library/view/hands-on-machine-learning/9781492032632/.}}{\nolinkurl{https://www.oreilly.com/library/view/hands-on-machine-learning/9781492032632/.}} This text is very useful since it contains many code examples and hands-on applications of all algorithms discussed in this course.
\end{enumerate}

\noindent
\end{frame}

\begin{frame}[plain,fragile]
\frametitle{Prerequisites}

Basic knowledge in programming and mathematics, with an emphasis on
linear algebra. Knowledge of Python or/and C++ as programming
languages is strongly recommended and experience with Jupiter notebook
is recommended. Required courses are the equivalents to the University
of Oslo mathematics courses MAT1100, MAT1110, MAT1120 and at least one
of the corresponding computing and programming courses INF1000/INF1110
or MAT-INF1100/MAT-INF1100L/BIOS1100/KJM-INF1100. Most universities
offer nowadays a basic programming course (often compulsory) where
Python is the recurring programming language.
\end{frame}

\begin{frame}[plain,fragile]
\frametitle{Learning outcomes}

\begin{block}{}

This course aims at giving you insights and knowledge about many of
the central algorithms used in Data Analysis and Machine Learning.
The course is project based and through various numerical projects,
normally three, you will be exposed to fundamental research problems
in these fields, with the aim to reproduce state of the art scientific
results. Both supervised and unsupervised methods will be covered. The
emphasis is on a frequentist approach, although we will try to link it
with a Bayesian approach as well. You will learn to develop and
structure large codes for studying different cases where Machine
Learning is applied to, get acquainted with computing facilities and
learn to handle large scientific projects. A good scientific and
ethical conduct is emphasized throughout the course. More
specifically, after this course you will

\begin{itemize}
\item Learn about basic data analysis, statistical analysis, Bayesian statistics, Monte Carlo sampling, data optimization and machine learning;

\item Be capable of extending the acquired knowledge to other systems and cases;

\item Have an understanding of central algorithms used in data analysis and machine learning;

\item Understand linear methods for regression and classification, from ordinary least squares, via Lasso and Ridge to Logistic regression;

\item Learn about neural networks and deep  learning methods for supervised and unsupervised learning. Emphasis on feed forward neural networks, convolutional and recurrent neural networks; 

\item Learn about about decision trees, random forests, bagging and boosting methods;

\item Learn about support vector machines and kernel transformations;

\item Reduction of data sets, from PCA to clustering;

\item Autoencoders and Reinforcement Learning;

\item Work on numerical projects to illustrate the theory. The projects play a central role and you are expected to know modern programming languages like Python or C++ and/or Fortran (Fortran2003 or later) or Julia or other.
\end{itemize}

\noindent
\end{block}
\end{frame}

\begin{frame}[plain,fragile]
\frametitle{Topics covered in this course: Statistical analysis and optimization of data}

The course has two central parts

\begin{enumerate}
\item Statistical analysis and optimization of data

\item Machine learning
\end{enumerate}

\noindent
These topics will be scattered thorughout the course and may not  necessarily be taught separately. Rather, we will often take an approach (during the lectures and project/exercise sessions) where say elements from statistical data analysis are mixed with specific Machine Learning algorithms

\begin{block}{Statistical analysis and optimization of data }

We plan to cover the following topics:
\begin{itemize}
\item Basic concepts, expectation values, variance, covariance, correlation functions and errors;

\item Simpler models, binomial distribution, the Poisson distribution, simple and multivariate normal distributions;

\item Central elements of Bayesian statistics and modeling;

\item Gradient methods for data optimization;

\item Monte Carlo methods, Markov chains, Gibbs sampling and Metropolis-Hastings sampling (tentative);

\item Estimation of errors and resampling techniques such as the cross-validation, blocking, bootstrapping and jackknife methods;

\item Principal Component Analysis (PCA) and its mathematical foundation;
\end{itemize}

\noindent
\end{block}
\end{frame}

\begin{frame}[plain,fragile]
\frametitle{Topics covered in this course}

\begin{block}{}
The following topics will be covered

\begin{itemize}
\item Pre deep-learning revolution (2008 approx)
\begin{itemize}

  \item Linear Regression and Logistic Regression, classification and regression problems;

  \item Bayesian linear and logistic regression, kernel regression;

  \item Decisions trees, Random Forests, Bagging and Boosting methods;

  \item Support vector machines (only survey);

  \item Unsupervised learning and  dimensionality reduction, from PCA to clustering;  

\end{itemize}

\noindent
\item Deep learning  
\begin{itemize}

  \item Neural networks and deep learning;

  \item Convolutional neural networks;

  \item Recurrent neural networks;

  \item Autoencoders

  \item Generative methods with an emphasis on Boltzmann Machines, Variational Autoencoders and Generalized Adversarial Networks;
\end{itemize}

\noindent
\end{itemize}

\noindent
Hands-on demonstrations, exercises and projects aim at deepening your understanding of these topics.

\end{block}
\end{frame}

\begin{frame}[plain,fragile]
\frametitle{Extremely useful tools, strongly recommended}

\begin{block}{and discussed at the lab sessions }
\begin{itemize}
  \item GIT for version control, and GitHub or GitLab as repositories, highly recommended. This will be discussed during the first exercise session

  \item Anaconda and other Python environments, see intro slides and links to programming resources at \href{{https://computationalscienceuio.github.io/RefreshProgrammingSkills/intro.html}}{\nolinkurl{https://computationalscienceuio.github.io/RefreshProgrammingSkills/intro.html}}
\end{itemize}

\noindent
\end{block}
\end{frame}

\begin{frame}[plain,fragile]
\frametitle{Other courses on Data science and Machine Learning  at UiO}

The link here \href{{https://www.mn.uio.no/english/research/about/centre-focus/innovation/data-science/studies/}}{\nolinkurl{https://www.mn.uio.no/english/research/about/centre-focus/innovation/data-science/studies/}}  gives an excellent overview of courses on Machine learning at UiO.
\begin{enumerate}
\item \href{{https://www.uio.no/studier/emner/matnat/fys/FYS5419/index-eng.html}}{FYS5419 Quantum Computing and Quantum Machine Learning}

\item \href{{https://www.uio.no/studier/emner/matnat/fys/FYS5429/index-eng.html}}{FYS5429 Advanced Machine Learning for the Physical Sciences}

\item \href{{http://www.uio.no/studier/emner/matnat/math/STK2100/index-eng.html}}{STK2100 Machine learning and statistical methods for prediction and classification}. 

\item \href{{https://www.uio.no/studier/emner/matnat/ifi/IN3050/index-eng.html}}{IN3050/4050 Introduction to Artificial Intelligence and Machine Learning}. Introductory course in machine learning and AI with an algorithmic approach. 

\item \href{{http://www.uio.no/studier/emner/matnat/math/STK-INF3000/index-eng.html}}{STK-INF3000/4000 Selected Topics in Data Science}. The course provides insight into selected contemporary relevant topics within Data Science. 

\item \href{{https://www.uio.no/studier/emner/matnat/ifi/IN4080/index.html}}{IN4080 Natural Language Processing}. Probabilistic and machine learning techniques applied to natural language processing. 

\item \href{{https://www.uio.no/studier/emner/matnat/math/STK-IN4300/index-eng.html}}{STK-IN4300 Statistical learning methods in Data Science}. An advanced introduction to statistical and machine learning. For students with a good mathematics and statistics background.

\item \href{{https://www.uio.no/studier/emner/matnat/ifi/IN4310/index.html}}{IN3310/4310 Deep Learnig for Image Analysis}

\item \href{{https://www.uio.no/studier/emner/matnat/math/STK4051/index-eng.html}}{STK4051 Computational Statistics}

\item \href{{https://www.uio.no/studier/emner/matnat/math/STK4021/index-eng.html}}{STK4021 Applied Bayesian Analysis and Numerical Methods}
\end{enumerate}

\noindent
\end{frame}

\begin{frame}[plain,fragile]
\frametitle{Introduction}

Our emphasis throughout this series of lectures  
is on understanding the mathematical aspects of
different algorithms used in the fields of data analysis and machine learning. 

However, where possible we will emphasize the
importance of using available software. We start thus with a hands-on
and top-down approach to machine learning. The aim is thus to start with
relevant data or data we have produced 
and use these to introduce statistical data analysis
concepts and machine learning algorithms before we delve into the
algorithms themselves. The examples we will use in the beginning, start with simple
polynomials with random noise added. We will use the Python
software package \href{{http://scikit-learn.org/stable/}}{Scikit-Learn} and
introduce various machine learning algorithms to make fits of
the data and predictions. We move thereafter to more interesting
cases such as data from say experiments (below we will look at experimental nuclear binding energies as an example).
These are examples where we can easily set up the data and
then use machine learning algorithms included in for example
\textbf{Scikit-Learn}. 

These examples will serve us the purpose of getting
started. Furthermore, they allow us to catch more than two birds with
a stone. They will allow us to bring in some programming specific
topics and tools as well as showing the power of various Python 
libraries for machine learning and statistical data analysis.  

Here, we will mainly focus on two
specific Python packages for Machine Learning, Scikit-Learn and
Tensorflow (see below for links etc).  Moreover, the examples we
introduce will serve as inputs to many of our discussions later, as
well as allowing you to set up models and produce your own data and
get started with programming.
\end{frame}

\begin{frame}[plain,fragile]
\frametitle{What is Machine Learning?}

Statistics, data science and machine learning form important fields of
research in modern science.  They describe how to learn and make
predictions from data, as well as allowing us to extract important
correlations about physical process and the underlying laws of motion
in large data sets. The latter, big data sets, appear frequently in
essentially all disciplines, from the traditional Science, Technology,
Mathematics and Engineering fields to Life Science, Law, education
research, the Humanities and the Social Sciences. 

It has become more
and more common to see research projects on big data in for example
the Social Sciences where extracting patterns from complicated survey
data is one of many research directions.  Having a solid grasp of data
analysis and machine learning is thus becoming central to scientific
computing in many fields, and competences and skills within the fields
of machine learning and scientific computing are nowadays strongly
requested by many potential employers. The latter cannot be
overstated, familiarity with machine learning has almost become a
prerequisite for many of the most exciting employment opportunities,
whether they are in bioinformatics, life science, physics or finance,
in the private or the public sector. This author has had several
students or met students who have been hired recently based on their
skills and competences in scientific computing and data science, often
with marginal knowledge of machine learning.

Machine learning is a subfield of computer science, and is closely
related to computational statistics.  It evolved from the study of
pattern recognition in artificial intelligence (AI) research, and has
made contributions to AI tasks like computer vision, natural language
processing and speech recognition. Many of the methods we will study are also 
strongly rooted in basic mathematics and physics research. 

Ideally, machine learning represents the science of giving computers
the ability to learn without being explicitly programmed.  The idea is
that there exist generic algorithms which can be used to find patterns
in a broad class of data sets without having to write code
specifically for each problem. The algorithm will build its own logic
based on the data.  You should however always keep in mind that
machines and algorithms are to a large extent developed by humans. The
insights and knowledge we have about a specific system, play a central
role when we develop a specific machine learning algorithm. 

Machine learning is an extremely rich field, in spite of its young
age. The increases we have seen during the last three decades in
computational capabilities have been followed by developments of
methods and techniques for analyzing and handling large date sets,
relying heavily on statistics, computer science and mathematics.  The
field is rather new and developing rapidly. Popular software packages
written in Python for machine learning like
\href{{http://scikit-learn.org/stable/}}{Scikit-learn},
\href{{https://www.tensorflow.org/}}{Tensorflow},
\href{{http://pytorch.org/}}{PyTorch} and \href{{https://keras.io/}}{Keras}, all
freely available at their respective GitHub sites, encompass
communities of developers in the thousands or more. And the number of
code developers and contributors keeps increasing. Not all the
algorithms and methods can be given a rigorous mathematical
justification, opening up thereby large rooms for experimenting and
trial and error and thereby exciting new developments.  However, a
solid command of linear algebra, multivariate theory, probability
theory, statistical data analysis, understanding errors and Monte
Carlo methods are central elements in a proper understanding of many
of algorithms and methods we will discuss.
\end{frame}

\begin{frame}[plain,fragile]
\frametitle{Types of Machine Learning}

The approaches to machine learning are many, but are often split into
two main categories.  In \emph{supervised learning} we know the answer to a
problem, and let the computer deduce the logic behind it. On the other
hand, \emph{unsupervised learning} is a method for finding patterns and
relationship in data sets without any prior knowledge of the system.
Some authours also operate with a third category, namely
\emph{reinforcement learning}. This is a paradigm of learning inspired by
behavioral psychology, where learning is achieved by trial-and-error,
solely from rewards and punishment.

Another way to categorize machine learning tasks is to consider the
desired output of a system.  Some of the most common tasks are:

\begin{itemize}
  \item Classification: Outputs are divided into two or more classes. The goal is to   produce a model that assigns inputs into one of these classes. An example is to identify  digits based on pictures of hand-written ones. Classification is typically supervised learning.

  \item Regression: Finding a functional relationship between an input data set and a reference data set.   The goal is to construct a function that maps input data to continuous output values.

  \item Clustering: Data are divided into groups with certain common traits, without knowing the different groups beforehand.  It is thus a form of unsupervised learning.
\end{itemize}

\noindent
\end{frame}

\begin{frame}[plain,fragile]
\frametitle{Essential elements of ML}

The methods we cover have three main topics in common, irrespective of
whether we deal with supervised or unsupervised learning.
\begin{itemize}
\pause
\item The first ingredient is normally our data set (which can be subdivided into training, validation  and test data). Many find the most difficult part of using Machine Learning to be the set up of your data in a meaningful way. 

\pause
\item The second item is a model which is normally a function of some parameters.  The model reflects our knowledge of the system (or lack thereof). As an example, if we know that our data show a behavior similar to what would be predicted by a polynomial, fitting our data to a polynomial of some degree would then determin our model. 

\pause
\item The last ingredient is a so-called \textbf{cost/loss} function (or error or risk function) which allows us to present an estimate on how good our model is in reproducing the data it is supposed to train.  
\end{itemize}

\noindent
\end{frame}

\begin{frame}[plain,fragile]
\frametitle{An optimization/minimization problem}

At the heart of basically all Machine Learning algorithms we will encounter so-called minimization or optimization algorithms. A large family of such methods are so-called \textbf{gradient methods}.
\end{frame}

\begin{frame}[plain,fragile]
\frametitle{A Frequentist approach to data analysis}

When you hear phrases like \textbf{predictions and estimations} and
\textbf{correlations and causations}, what do you think of?  May be you think
of the difference between classifying new data points and generating
new data points.
Or perhaps you consider that correlations represent some kind of symmetric statements like
if $A$ is correlated with $B$, then $B$ is correlated with
$A$. Causation on the other hand is directional, that is if $A$ causes $B$, $B$ does not
necessarily cause $A$.

These concepts are in some sense the difference between machine
learning and statistics. In machine learning and prediction based
tasks, we are often interested in developing algorithms that are
capable of learning patterns from given data in an automated fashion,
and then using these learned patterns to make predictions or
assessments of newly given data. In many cases, our primary concern
is the quality of the predictions or assessments, and we are less
concerned about the underlying patterns that were learned in order
to make these predictions.

In machine learning we normally use \href{{https://en.wikipedia.org/wiki/Frequentist_inference}}{a so-called frequentist approach},
where the aim is to make predictions and find correlations. We focus
less on for example extracting a probability distribution function (PDF). The PDF can be
used in turn to make estimations and find causations such as given $A$
what is the likelihood of finding $B$.
\end{frame}

\begin{frame}[plain,fragile]
\frametitle{What is a good model?}

In science and engineering we often end up in situations where we want to infer (or learn) a
quantitative model $M$ for a given set of sample points $\bm{X} \in [x_1, x_2,\dots x_N]$.

As we will see repeatedely in these lectures, we could try to fit these data points to a model given by a
straight line, or if we wish to be more sophisticated to a more complex
function.

The reason for inferring such a model is that it
serves many useful purposes. On the one hand, the model can reveal information
encoded in the data or underlying mechanisms from which the data were generated. For instance, we could discover important
corelations that relate interesting physics interpretations.

In addition, it can simplify the representation of the given data set and help
us in making predictions about  future data samples.

A first important consideration to keep in mind is that inferring the \emph{correct} model
for a given data set is an elusive, if not impossible, task. The fundamental difficulty
is that if we are not specific about what we mean by a \emph{correct} model, there
could easily be many different models that fit the given data set \emph{equally well}.
\end{frame}

\begin{frame}[plain,fragile]
\frametitle{What is a good model? Can we define it?}

The central question is this: what leads us to say that a model is correct or
optimal for a given data set? To make the model inference problem well posed, i.e.,
to guarantee that there is a unique optimal model for the given data, we need to
impose additional assumptions or restrictions on the class of models considered. To
this end, we should not be looking for just any model that can describe the data.
Instead, we should look for a \textbf{model} $M$ that is the best among a restricted class
of models. In addition, to make the model inference problem computationally
tractable, we need to specify how restricted the class of models needs to be. A
common strategy is to start 
with the simplest possible class of models that is just necessary to describe the data
or solve the problem at hand. More precisely, the model class should be rich enough
to contain at least one model that can fit the data to a desired accuracy and yet be
restricted enough that it is relatively simple to find the best model for the given data.

Thus, the most popular strategy is to start from the
simplest class of models and increase the complexity of the models only when the
simpler models become inadequate. For instance, if we work with a regression problem to fit a set of sample points, one
may first try the simplest class of models, namely linear models, followed obviously by more complex models.

How to evaluate which model fits best the data is something we will come back to over and over again in these sets of lectures.
\end{frame}

\begin{frame}[plain,fragile]
\frametitle{Software and needed installations}

We will make extensive use of Python as programming language and its
myriad of available libraries.  You will find
Jupyter notebooks invaluable in your work.  You can run \textbf{R}
codes in the Jupyter/IPython notebooks, with the immediate benefit of
visualizing your data. You can also use compiled languages like C++,
Rust, Julia, Fortran etc if you prefer. The focus in these lectures will be
on Python.

If you have Python installed (we strongly recommend Python3) and you feel
pretty familiar with installing different packages, we recommend that
you install the following Python packages via \textbf{pip} as 

\begin{enumerate}
\item pip install numpy scipy matplotlib ipython scikit-learn mglearn sympy pandas pillow 
\end{enumerate}

\noindent
For Python3, replace \textbf{pip} with \textbf{pip3}.

For OSX users we recommend, after having installed Xcode, to
install \textbf{brew}. Brew allows for a seamless installation of additional
software via for example 

\begin{enumerate}
\item brew install python3
\end{enumerate}

\noindent
For Linux users, with its variety of distributions like for example the widely popular Ubuntu distribution,
you can use \textbf{pip} as well and simply install Python as 

\begin{enumerate}
\item sudo apt-get install python3  (or python for pyhton2.7)
\end{enumerate}

\noindent
etc etc.
\end{frame}

\begin{frame}[plain,fragile]
\frametitle{Python installers}

If you don't want to perform these operations separately and venture
into the hassle of exploring how to set up dependencies and paths, we
recommend two widely used distrubutions which set up all relevant
dependencies for Python, namely 

\begin{itemize}
\item \href{{https://docs.anaconda.com/}}{Anaconda}, 
\end{itemize}

\noindent
which is an open source
distribution of the Python and R programming languages for large-scale
data processing, predictive analytics, and scientific computing, that
aims to simplify package management and deployment. Package versions
are managed by the package management system \textbf{conda}. 

\begin{itemize}
\item \href{{https://www.enthought.com/product/canopy/}}{Enthought canopy} 
\end{itemize}

\noindent
is a Python
distribution for scientific and analytic computing distribution and
analysis environment, available for free and under a commercial
license.

Furthermore, \href{{https://colab.research.google.com/notebooks/welcome.ipynb}}{Google's Colab} is a free Jupyter notebook environment that requires 
no setup and runs entirely in the cloud. Try it out!
\end{frame}

\begin{frame}[plain,fragile]
\frametitle{Useful Python libraries}

Here we list several useful Python libraries we strongly recommend (if you use anaconda many of these are already there)

\begin{itemize}
\item \href{{https://www.numpy.org/}}{NumPy} is a highly popular library for large, multi-dimensional arrays and matrices, along with a large collection of high-level mathematical functions to operate on these arrays

\item \href{{https://pandas.pydata.org/}}{The pandas} library provides high-performance, easy-to-use data structures and data analysis tools 

\item \href{{http://xarray.pydata.org/en/stable/}}{Xarray} is a Python package that makes working with labelled multi-dimensional arrays simple, efficient, and fun!

\item \href{{https://www.scipy.org/}}{Scipy} (pronounced “Sigh Pie”) is a Python-based ecosystem of open-source software for mathematics, science, and engineering. 

\item \href{{https://matplotlib.org/}}{Matplotlib} is a Python 2D plotting library which produces publication quality figures in a variety of hardcopy formats and interactive environments across platforms.

\item \href{{https://github.com/HIPS/autograd}}{Autograd} can automatically differentiate native Python and Numpy code. It can handle a large subset of Python's features, including loops, ifs, recursion and closures, and it can even take derivatives of derivatives of derivatives

\item \href{{https://jax.readthedocs.io/en/latest/index.html}}{JAX} has now more or less replaced \textbf{Autograd}. JAX is Autograd and XLA, brought together for high-performance numerical computing and machine learning research. It provides composable transformations of Python+NumPy programs: differentiate, vectorize, parallelize, Just-In-Time compile to GPU/TPU, and more.

\item \href{{https://www.sympy.org/en/index.html}}{SymPy} is a Python library for symbolic mathematics. 

\item \href{{https://scikit-learn.org/stable/}}{scikit-learn} has simple and efficient tools for machine learning, data mining and data analysis

\item \href{{https://www.tensorflow.org/}}{TensorFlow} is a Python library for fast numerical computing created and released by Google

\item \href{{https://keras.io/}}{Keras} is a high-level neural networks API, written in Python and capable of running on top of TensorFlow, CNTK, or Theano

\item And many more such as \href{{https://pytorch.org/}}{pytorch},  \href{{https://pypi.org/project/Theano/}}{Theano} etc 
\end{itemize}

\noindent
\end{frame}

\begin{frame}[plain,fragile]
\frametitle{Installing R, C++, cython or Julia}

You will also find it convenient to utilize \textbf{R}. We will mainly
use Python during our lectures and in various projects and exercises.
Those of you
already familiar with \textbf{R} should feel free to continue using \textbf{R}, keeping
however an eye on the parallel Python set ups. Similarly, if you are a
Python afecionado, feel free to explore \textbf{R} as well.  Jupyter/Ipython
notebook allows you to run \textbf{R} codes interactively in your
browser. The software library \textbf{R} is really tailored  for statistical data analysis
and allows for an easy usage of the tools and algorithms we will discuss in these
lectures.

To install \textbf{R} with Jupyter notebook 
\href{{https://mpacer.org/maths/r-kernel-for-ipython-notebook}}{follow the link here}
\end{frame}

\begin{frame}[plain,fragile]
\frametitle{Installing R, C++, cython, Numba etc}

For the C++ aficionados, Jupyter/IPython notebook allows you also to
install C++ and run codes written in this language interactively in
the browser. Since we will emphasize writing many of the algorithms
yourself, you can thus opt for either Python or C++ (or Fortran or other compiled languages) as programming
languages.

To add more entropy, \textbf{cython} can also be used when running your
notebooks. It means that Python with the jupyter notebook
setup allows you to integrate widely popular softwares and tools for
scientific computing. Similarly, the 
\href{{https://numba.pydata.org/}}{Numba Python package} delivers increased performance
capabilities with minimal rewrites of your codes.  With its
versatility, including symbolic operations, Python offers a unique
computational environment. Your jupyter notebook can easily be
converted into a nicely rendered \textbf{PDF} file or a Latex file for
further processing. For example, convert to latex as 



\begin{minted}[fontsize=\fontsize{9pt}{9pt},linenos=false,mathescape,baselinestretch=1.0,fontfamily=tt,xleftmargin=2mm]{text}
pycod jupyter nbconvert filename.ipynb --to latex 

\end{minted}


And to add more versatility, the Python package \href{{http://www.sympy.org/en/index.html}}{SymPy} is a Python library for symbolic mathematics. It aims to become a full-featured computer algebra system (CAS)  and is entirely written in Python. 

Finally, if you wish to use the light mark-up language 
\href{{https://github.com/hplgit/doconce}}{doconce} you can convert a standard ascii text file into various HTML 
formats, ipython notebooks, latex files, pdf files etc with minimal edits. These lectures were generated using \textbf{doconce}.
\end{frame}

\begin{frame}[plain,fragile]
\frametitle{Numpy examples and Important Matrix and vector handling packages}

There are several central software libraries for linear algebra and eigenvalue problems. Several of the more
popular ones have been wrapped into ofter software packages like those from the widely used text \textbf{Numerical Recipes}. The original source codes in many of the available packages are often taken from the widely used
software package LAPACK, which follows two other popular packages
developed in the 1970s, namely EISPACK and LINPACK.  We describe them shortly here.

\begin{itemize}
  \item LINPACK: package for linear equations and least square problems.

  \item LAPACK:package for solving symmetric, unsymmetric and generalized eigenvalue problems. From LAPACK's website \href{{http://www.netlib.org}}{\nolinkurl{http://www.netlib.org}} it is possible to download for free all source codes from this library. Both C/C++ and Fortran versions are available.

  \item BLAS (I, II and III): (Basic Linear Algebra Subprograms) are routines that provide standard building blocks for performing basic vector and matrix operations. Blas I is vector operations, II vector-matrix operations and III matrix-matrix operations. Highly parallelized and efficient codes, all available for download from \href{{http://www.netlib.org}}{\nolinkurl{http://www.netlib.org}}.
\end{itemize}

\noindent
\end{frame}

\begin{frame}[plain,fragile]
\frametitle{Basic Matrix Features}

\begin{block}{Matrix properties reminder }
\[
 \mathbf{A} =
      \begin{bmatrix} a_{11} & a_{12} & a_{13} & a_{14} \\
                                 a_{21} & a_{22} & a_{23} & a_{24} \\
                                   a_{31} & a_{32} & a_{33} & a_{34} \\
                                  a_{41} & a_{42} & a_{43} & a_{44}
             \end{bmatrix}\qquad
\mathbf{I} =
      \begin{bmatrix} 1 & 0 & 0 & 0 \\
                                 0 & 1 & 0 & 0 \\
                                 0 & 0 & 1 & 0 \\
                                 0 & 0 & 0 & 1
             \end{bmatrix}
\]

The inverse of a matrix is defined by

\[
\mathbf{A}^{-1} \cdot \mathbf{A} = I
\]


{\footnotesize
\begin{tabular}{ccc}
\hline
\multicolumn{1}{c}{ Relations } & \multicolumn{1}{c}{ Name } & \multicolumn{1}{c}{ matrix elements } \\
\hline
$A=A^{T}$                          & symmetric       & $a_{ij}=a_{ji}$                                            \\
$A=\left (A^{T}\right )^{-1}$      & real orthogonal & $\sum_k a_{ik}a_{jk}=\sum_k a_{ki} a_{kj}=\delta_{ij}$     \\
$A=A^*$                            & real matrix     & $a_{ij}=a_{ij}^*$                                          \\
$A=A^{\dagger}$                    & hermitian       & $a_{ij}=a_{ji}^*$                                          \\
$A=\left(A^{\dagger}\right )^{-1}$ & unitary         & $\sum_k a_{ik}a_{jk}^*=\sum_k a_{ki}^* a_{kj}=\delta_{ij}$ \\
\hline
\end{tabular}
}

\noindent
\end{block}
\end{frame}

\begin{frame}[plain,fragile]
% No title on this slide

\noindent\textbf{Some famous Matrices.}
\begin{itemize}
  \item Diagonal if $a_{ij}=0$ for $i\ne j$

  \item Upper triangular if $a_{ij}=0$ for $i>j$

  \item Lower triangular if $a_{ij}=0$ for $i<j$

  \item Upper Hessenberg if $a_{ij}=0$ for $i>j+1$

  \item Lower Hessenberg if $a_{ij}=0$ for $i<j+1$

  \item Tridiagonal if $a_{ij}=0$ for $|i -j|>1$

  \item Lower banded with bandwidth $p$: $a_{ij}=0$ for $i>j+p$

  \item Upper banded with bandwidth $p$: $a_{ij}=0$ for $i<j+p$

  \item Banded, block upper triangular, block lower triangular....
\end{itemize}

\noindent
\end{frame}

\begin{frame}[plain,fragile]
% No title on this slide

\noindent\textbf{More Basic Matrix Features.}
\begin{block}{Some Equivalent Statements }
For an $N\times N$ matrix  $\mathbf{A}$ the following properties are all equivalent

\begin{itemize}
  \item If the inverse of $\mathbf{A}$ exists, $\mathbf{A}$ is nonsingular.

  \item The equation $\mathbf{Ax}=0$ implies $\mathbf{x}=0$.

  \item The rows of $\mathbf{A}$ form a basis of $R^N$.

  \item The columns of $\mathbf{A}$ form a basis of $R^N$.

  \item $\mathbf{A}$ is a product of elementary matrices.

  \item $0$ is not eigenvalue of $\mathbf{A}$.
\end{itemize}

\noindent
\end{block}
\end{frame}

\begin{frame}[plain,fragile]
\frametitle{Numpy and arrays}

\href{{http://www.numpy.org/}}{Numpy} provides an easy way to handle arrays in Python. The standard way to import this library is as



\begin{minted}[fontsize=\fontsize{9pt}{9pt},linenos=false,mathescape,baselinestretch=1.0,fontfamily=tt,xleftmargin=2mm]{python}
import numpy as np

\end{minted}

Here follows a simple example where we set up an array of ten elements, all determined by random numbers drawn according to the normal distribution,




\begin{minted}[fontsize=\fontsize{9pt}{9pt},linenos=false,mathescape,baselinestretch=1.0,fontfamily=tt,xleftmargin=2mm]{python}
n = 10
x = np.random.normal(size=n)
print(x)

\end{minted}

We defined a vector $x$ with $n=10$ elements with its values given by the Normal distribution $N(0,1)$.
Another alternative is to declare a vector as follows




\begin{minted}[fontsize=\fontsize{9pt}{9pt},linenos=false,mathescape,baselinestretch=1.0,fontfamily=tt,xleftmargin=2mm]{python}
import numpy as np
x = np.array([1, 2, 3])
print(x)

\end{minted}

Here we have defined a vector with three elements, with $x_0=1$, $x_1=2$ and $x_2=3$. Note that both Python and C++
start numbering array elements from $0$ and on. This means that a vector with $n$ elements has a sequence of entities $x_0, x_1, x_2, \dots, x_{n-1}$. We could also let (recommended) Numpy to compute the logarithms of a specific array as




\begin{minted}[fontsize=\fontsize{9pt}{9pt},linenos=false,mathescape,baselinestretch=1.0,fontfamily=tt,xleftmargin=2mm]{python}
import numpy as np
x = np.log(np.array([4, 7, 8]))
print(x)

\end{minted}


In the last example we used Numpy's unary function $np.log$. This function is
highly tuned to compute array elements since the code is vectorized
and does not require looping. We normaly recommend that you use the
Numpy intrinsic functions instead of the corresponding \textbf{log} function
from Python's \textbf{math} module. The looping is done explicitely by the
\textbf{np.log} function. The alternative, and slower way to compute the
logarithms of a vector would be to write








\begin{minted}[fontsize=\fontsize{9pt}{9pt},linenos=false,mathescape,baselinestretch=1.0,fontfamily=tt,xleftmargin=2mm]{python}
import numpy as np
from math import log
x = np.array([4, 7, 8])
for i in range(0, len(x)):
    x[i] = log(x[i])
print(x)

\end{minted}

We note that our code is much longer already and we need to import the \textbf{log} function from the \textbf{math} module. 
The attentive reader will also notice that the output is $[1, 1, 2]$. Python interprets automagically our numbers as integers (like the \textbf{automatic} keyword in C++). To change this we could define our array elements to be double precision numbers as




\begin{minted}[fontsize=\fontsize{9pt}{9pt},linenos=false,mathescape,baselinestretch=1.0,fontfamily=tt,xleftmargin=2mm]{python}
import numpy as np
x = np.log(np.array([4, 7, 8], dtype = np.float64))
print(x)

\end{minted}

or simply write them as double precision numbers (Python uses 64 bits as default for floating point type variables), that is




\begin{minted}[fontsize=\fontsize{9pt}{9pt},linenos=false,mathescape,baselinestretch=1.0,fontfamily=tt,xleftmargin=2mm]{python}
import numpy as np
x = np.log(np.array([4.0, 7.0, 8.0]))
print(x)

\end{minted}

To check the number of bytes (remember that one byte contains eight bits for double precision variables), you can use simple use the \textbf{itemsize} functionality (the array $x$ is actually an object which inherits the functionalities defined in Numpy) as 




\begin{minted}[fontsize=\fontsize{9pt}{9pt},linenos=false,mathescape,baselinestretch=1.0,fontfamily=tt,xleftmargin=2mm]{python}
import numpy as np
x = np.log(np.array([4.0, 7.0, 8.0]))
print(x.itemsize)

\end{minted}
\end{frame}

\begin{frame}[plain,fragile]
\frametitle{Matrices in Python}

Having defined vectors, we are now ready to try out matrices. We can
define a $3 \times 3 $ real matrix $\bm{A}$ as (recall that we user
lowercase letters for vectors and uppercase letters for matrices)





\begin{minted}[fontsize=\fontsize{9pt}{9pt},linenos=false,mathescape,baselinestretch=1.0,fontfamily=tt,xleftmargin=2mm]{python}
import numpy as np
A = np.log(np.array([ [4.0, 7.0, 8.0], [3.0, 10.0, 11.0], [4.0, 5.0, 7.0] ]))
print(A)

\end{minted}

If we use the \textbf{shape} function we would get $(3, 3)$ as output, that is verifying that our matrix is a $3\times 3$ matrix. We can slice the matrix and print for example the first column (Python organized matrix elements in a row-major order, see below) as





\begin{minted}[fontsize=\fontsize{9pt}{9pt},linenos=false,mathescape,baselinestretch=1.0,fontfamily=tt,xleftmargin=2mm]{python}
import numpy as np
A = np.log(np.array([ [4.0, 7.0, 8.0], [3.0, 10.0, 11.0], [4.0, 5.0, 7.0] ]))
# print the first column, row-major order and elements start with 0
print(A[:,0]) 

\end{minted}

We can continue this was by printing out other columns or rows. The example here prints out the second column





\begin{minted}[fontsize=\fontsize{9pt}{9pt},linenos=false,mathescape,baselinestretch=1.0,fontfamily=tt,xleftmargin=2mm]{python}
import numpy as np
A = np.log(np.array([ [4.0, 7.0, 8.0], [3.0, 10.0, 11.0], [4.0, 5.0, 7.0] ]))
# print the first column, row-major order and elements start with 0
print(A[1,:]) 

\end{minted}

Numpy contains many other functionalities that allow us to slice, subdivide etc etc arrays. We strongly recommend that you look up the \href{{http://www.numpy.org/}}{Numpy website for more details}. Useful functions when defining a matrix are the \textbf{np.zeros} function which declares a matrix of a given dimension and sets all elements to zero






\begin{minted}[fontsize=\fontsize{9pt}{9pt},linenos=false,mathescape,baselinestretch=1.0,fontfamily=tt,xleftmargin=2mm]{python}
import numpy as np
n = 10
# define a matrix of dimension 10 x 10 and set all elements to zero
A = np.zeros( (n, n) )
print(A) 

\end{minted}

or initializing all elements to 






\begin{minted}[fontsize=\fontsize{9pt}{9pt},linenos=false,mathescape,baselinestretch=1.0,fontfamily=tt,xleftmargin=2mm]{python}
import numpy as np
n = 10
# define a matrix of dimension 10 x 10 and set all elements to one
A = np.ones( (n, n) )
print(A) 

\end{minted}

or as unitarily distributed random numbers (see the material on random number generators in the statistics part)






\begin{minted}[fontsize=\fontsize{9pt}{9pt},linenos=false,mathescape,baselinestretch=1.0,fontfamily=tt,xleftmargin=2mm]{python}
import numpy as np
n = 10
# define a matrix of dimension 10 x 10 and set all elements to random numbers with x \in [0, 1]
A = np.random.rand(n, n)
print(A) 

\end{minted}


As we will see throughout these lectures, there are several extremely useful functionalities in Numpy.
As an example, consider the discussion of the covariance matrix. Suppose we have defined three vectors
$\bm{x}, \bm{y}, \bm{z}$ with $n$ elements each. The covariance matrix is defined as 
\[
\bm{\Sigma} = \begin{bmatrix} \sigma_{xx} & \sigma_{xy} & \sigma_{xz} \\
                              \sigma_{yx} & \sigma_{yy} & \sigma_{yz} \\
                              \sigma_{zx} & \sigma_{zy} & \sigma_{zz} 
             \end{bmatrix},
\]
where for example
\[
\sigma_{xy} =\frac{1}{n} \sum_{i=0}^{n-1}(x_i- \overline{x})(y_i- \overline{y}).
\]
The Numpy function \textbf{np.cov} calculates the covariance elements using the factor $1/(n-1)$ instead of $1/n$ since it assumes we do not have the exact mean values. 
The following simple function uses the \textbf{np.vstack} function which takes each vector of dimension $1\times n$ and produces a $3\times n$ matrix $\bm{W}$
\[
\bm{W} = \begin{bmatrix} x_0 & x_1 & x_2 & \dots & x_{n-2} & x_{n-1} \\
                         y_0 & y_1 & y_2 & \dots & y_{n-2} & y_{n-1} \\
			 z_0 & z_1 & z_2 & \dots & z_{n-2} & z_{n-1} \\
             \end{bmatrix},
\]

which in turn is converted into into the $3\times 3$ covariance matrix
$\bm{\Sigma}$ via the Numpy function \textbf{np.cov()}. We note that we can also calculate
the mean value of each set of samples $\bm{x}$ etc using the Numpy
function \textbf{np.mean(x)}. We can also extract the eigenvalues of the
covariance matrix through the \textbf{np.linalg.eig()} function.

















\begin{minted}[fontsize=\fontsize{9pt}{9pt},linenos=false,mathescape,baselinestretch=1.0,fontfamily=tt,xleftmargin=2mm]{python}
# Importing various packages
import numpy as np

n = 100
x = np.random.normal(size=n)
print(np.mean(x))
y = 4+3*x+np.random.normal(size=n)
print(np.mean(y))
z = x**3+np.random.normal(size=n)
print(np.mean(z))
W = np.vstack((x, y, z))
Sigma = np.cov(W)
print(Sigma)
Eigvals, Eigvecs = np.linalg.eig(Sigma)
print(Eigvals)

\end{minted}














\begin{minted}[fontsize=\fontsize{9pt}{9pt},linenos=false,mathescape,baselinestretch=1.0,fontfamily=tt,xleftmargin=2mm]{python}
import numpy as np
import matplotlib.pyplot as plt
from scipy import sparse
eye = np.eye(4)
print(eye)
sparse_mtx = sparse.csr_matrix(eye)
print(sparse_mtx)
x = np.linspace(-10,10,100)
y = np.sin(x)
plt.plot(x,y,marker='x')
plt.show()

\end{minted}
\end{frame}

\begin{frame}[plain,fragile]
\frametitle{Meet the Pandas}

\vspace{6mm}

% inline figure
\centerline{\includegraphics[width=0.8\linewidth]{fig/pandas.jpg}}

\vspace{6mm}

Another useful Python package is
\href{{https://pandas.pydata.org/}}{pandas}, which is an open source library
providing high-performance, easy-to-use data structures and data
analysis tools for Python. \textbf{pandas} stands for panel data, a term borrowed from econometrics and is an efficient library for data analysis with an emphasis on tabular data.
\textbf{pandas} has two major classes, the \textbf{DataFrame} class with two-dimensional data objects and tabular data organized in columns and the class \textbf{Series} with a focus on one-dimensional data objects. Both classes allow you to index data easily as we will see in the examples below. 
\textbf{pandas} allows you also to perform mathematical operations on the data, spanning from simple reshapings of vectors and matrices to statistical operations. 

The following simple example shows how we can, in an easy way make tables of our data. Here we define a data set which includes names, place of birth and date of birth, and displays the data in an easy to read way. We will see repeated use of \textbf{pandas}, in particular in connection with classification of data. 











\begin{minted}[fontsize=\fontsize{9pt}{9pt},linenos=false,mathescape,baselinestretch=1.0,fontfamily=tt,xleftmargin=2mm]{python}
import pandas as pd
from IPython.display import display
data = {'First Name': ["Frodo", "Bilbo", "Aragorn II", "Samwise"],
        'Last Name': ["Baggins", "Baggins","Elessar","Gamgee"],
        'Place of birth': ["Shire", "Shire", "Eriador", "Shire"],
        'Date of Birth T.A.': [2968, 2890, 2931, 2980]
        }
data_pandas = pd.DataFrame(data)
display(data_pandas)

\end{minted}


In the above we have imported \textbf{pandas} with the shorthand \textbf{pd}, the latter has become the standard way we import \textbf{pandas}. We make then a list of various variables
and reorganize the aboves lists into a \textbf{DataFrame} and then print out  a neat table with specific column labels as \emph{Name}, \emph{place of birth} and \emph{date of birth}.
Displaying these results, we see that the indices are given by the default numbers from zero to three.
\textbf{pandas} is extremely flexible and we can easily change the above indices by defining a new type of indexing as



\begin{minted}[fontsize=\fontsize{9pt}{9pt},linenos=false,mathescape,baselinestretch=1.0,fontfamily=tt,xleftmargin=2mm]{python}
data_pandas = pd.DataFrame(data,index=['Frodo','Bilbo','Aragorn','Sam'])
display(data_pandas)

\end{minted}

Thereafter we display the content of the row which begins with the index \textbf{Aragorn}


\begin{minted}[fontsize=\fontsize{9pt}{9pt},linenos=false,mathescape,baselinestretch=1.0,fontfamily=tt,xleftmargin=2mm]{python}
display(data_pandas.loc['Aragorn'])

\end{minted}


We can easily append data to this, for example








\begin{minted}[fontsize=\fontsize{9pt}{9pt},linenos=false,mathescape,baselinestretch=1.0,fontfamily=tt,xleftmargin=2mm]{python}
new_hobbit = {'First Name': ["Peregrin"],
              'Last Name': ["Took"],
              'Place of birth': ["Shire"],
              'Date of Birth T.A.': [2990]
              }
data_pandas=data_pandas.append(pd.DataFrame(new_hobbit, index=['Pippin']))
display(data_pandas)

\end{minted}


Here are other examples where we use the \textbf{DataFrame} functionality to handle arrays, now with more interesting features for us, namely numbers. We set up a matrix 
of dimensionality $10\times 5$ and compute the mean value and standard deviation of each column. Similarly, we can perform mathematial operations like squaring the matrix elements and many other operations. 














\begin{minted}[fontsize=\fontsize{9pt}{9pt},linenos=false,mathescape,baselinestretch=1.0,fontfamily=tt,xleftmargin=2mm]{python}
import numpy as np
import pandas as pd
from IPython.display import display
np.random.seed(100)
# setting up a 10 x 5 matrix
rows = 10
cols = 5
a = np.random.randn(rows,cols)
df = pd.DataFrame(a)
display(df)
print(df.mean())
print(df.std())
display(df**2)

\end{minted}


Thereafter we can select specific columns only and plot final results



















\begin{minted}[fontsize=\fontsize{9pt}{9pt},linenos=false,mathescape,baselinestretch=1.0,fontfamily=tt,xleftmargin=2mm]{python}
df.columns = ['First', 'Second', 'Third', 'Fourth', 'Fifth']
df.index = np.arange(10)

display(df)
print(df['Second'].mean() )
print(df.info())
print(df.describe())

from pylab import plt, mpl
plt.style.use('seaborn')
mpl.rcParams['font.family'] = 'serif'

df.cumsum().plot(lw=2.0, figsize=(10,6))
plt.show()


df.plot.bar(figsize=(10,6), rot=15)
plt.show()

\end{minted}

We can produce a $4\times 4$ matrix





\begin{minted}[fontsize=\fontsize{9pt}{9pt},linenos=false,mathescape,baselinestretch=1.0,fontfamily=tt,xleftmargin=2mm]{python}
b = np.arange(16).reshape((4,4))
print(b)
df1 = pd.DataFrame(b)
print(df1)

\end{minted}

and many other operations. 

The \textbf{Series} class is another important class included in
\textbf{pandas}. You can view it as a specialization of \textbf{DataFrame} but where
we have just a single column of data. It shares many of the same features as \textbf{DataFrame}. As with \textbf{DataFrame},
most operations are vectorized, achieving thereby a high performance when dealing with computations of arrays, in particular labeled arrays.
As we will see below it leads also to a very concice code close to the mathematical operations we may be interested in.
For multidimensional arrays, we recommend strongly \href{{http://xarray.pydata.org/en/stable/}}{xarray}. \textbf{xarray} has much of the same flexibility as \textbf{pandas}, but allows for the extension to higher dimensions than two. We will see examples later of the usage of both \textbf{pandas} and \textbf{xarray}.
\end{frame}

\begin{frame}[plain,fragile]
\frametitle{Friday August 27}

"Video of Lecture August 27, 2021":"https://www.uio.no/studier/emner/matnat/fys/FYS-STK4155/h21/forelesningsvideoer/LectureThursdayAugust27.mp4?vrtx=view-as-webpage

\href{{https://www.uio.no/studier/emner/matnat/fys/FYS-STK3155/h20/forelesningsvideoer/LectureAug21.mp4?vrtx=view-as-webpage}}{Video of Lecture from fall 2020}  and \href{{https://github.com/CompPhysics/MachineLearning/blob/master/doc/HandWrittenNotes/NotesAugust21.pdf}}{Handwritten notes}
\end{frame}

\begin{frame}[plain,fragile]
\frametitle{A first summary}

\noindent\textbf{Simple linear regression model using \textbf{scikit-learn}.}
We start with perhaps our simplest possible example, using \textbf{Scikit-Learn} to perform linear regression analysis on a data set produced by us. 

What follows is a simple Python code where we have defined a function
$y$ in terms of the variable $x$. Both are defined as vectors with  $100$ entries. 
The numbers in the vector $\bm{x}$ are given
by random numbers generated with a uniform distribution with entries
$x_i \in [0,1]$ (more about probability distribution functions
later). These values are then used to define a function $y(x)$
(tabulated again as a vector) with a linear dependence on $x$ plus a
random noise added via the normal distribution.

The Numpy functions are imported used the \textbf{import numpy as np}
statement and the random number generator for the uniform distribution
is called using the function \textbf{np.random.rand()}, where we specificy
that we want $100$ random variables.  Using Numpy we define
automatically an array with the specified number of elements, $100$ in
our case.  With the Numpy function \textbf{randn()} we can compute random
numbers with the normal distribution (mean value $\mu$ equal to zero and
variance $\sigma^2$ set to one) and produce the values of $y$ assuming a linear
dependence as function of $x$

\[
y = 2x+N(0,1),
\]

where $N(0,1)$ represents random numbers generated by the normal
distribution.  From \textbf{Scikit-Learn} we import then the
\textbf{LinearRegression} functionality and make a prediction $\tilde{y} =
\alpha + \beta x$ using the function \textbf{fit(x,y)}. We call the set of
data $(\bm{x},\bm{y})$ for our training data. The Python package
\textbf{scikit-learn} has also a functionality which extracts the above
fitting parameters $\alpha$ and $\beta$ (see below). Later we will
distinguish between training data and test data.

For plotting we use the Python package
\href{{https://matplotlib.org/}}{matplotlib} which produces publication
quality figures. Feel free to explore the extensive
\href{{https://matplotlib.org/gallery/index.html}}{gallery} of examples. In
this example we plot our original values of $x$ and $y$ as well as the
prediction \textbf{ypredict} ($\tilde{y}$), which attempts at fitting our
data with a straight line.

The Python code follows here.




















\begin{minted}[fontsize=\fontsize{9pt}{9pt},linenos=false,mathescape,baselinestretch=1.0,fontfamily=tt,xleftmargin=2mm]{python}
# Importing various packages
import numpy as np
import matplotlib.pyplot as plt
from sklearn.linear_model import LinearRegression

x = np.random.rand(100,1)
y = 2*x+np.random.randn(100,1)
linreg = LinearRegression()
linreg.fit(x,y)
xnew = np.array([[0],[1]])
ypredict = linreg.predict(xnew)

plt.plot(xnew, ypredict, "r-")
plt.plot(x, y ,'ro')
plt.axis([0,1.0,0, 5.0])
plt.xlabel(r'$x$')
plt.ylabel(r'$y$')
plt.title(r'Simple Linear Regression')
plt.show()

\end{minted}


This example serves several aims. It allows us to demonstrate several
aspects of data analysis and later machine learning algorithms. The
immediate visualization shows that our linear fit is not
impressive. It goes through the data points, but there are many
outliers which are not reproduced by our linear regression.  We could
now play around with this small program and change for example the
factor in front of $x$ and the normal distribution.  Try to change the
function $y$ to

\[
y = 10x+0.01 \times N(0,1),
\]

where $x$ is defined as before.  Does the fit look better? Indeed, by
reducing the role of the noise given by the normal distribution we see immediately that
our linear prediction seemingly reproduces better the training
set. However, this testing 'by the eye' is obviouly not satisfactory in the
long run. Here we have only defined the training data and our model, and 
have not discussed a more rigorous approach to the \textbf{cost} function.

We need more rigorous criteria in defining whether we have succeeded or
not in modeling our training data.  You will be surprised to see that
many scientists seldomly venture beyond this 'by the eye' approach. A
standard approach for the \emph{cost} function is the so-called $\chi^2$
function (a variant of the mean-squared error (MSE))

\[ \chi^2 = \frac{1}{n}
\sum_{i=0}^{n-1}\frac{(y_i-\tilde{y}_i)^2}{\sigma_i^2}, 
\] 

where $\sigma_i^2$ is the variance (to be defined later) of the entry
$y_i$.  We may not know the explicit value of $\sigma_i^2$, it serves
however the aim of scaling the equations and make the cost function
dimensionless.  

Minimizing the cost function is a central aspect of
our discussions to come. Finding its minima as function of the model
parameters ($\alpha$ and $\beta$ in our case) will be a recurring
theme in these series of lectures. Essentially all machine learning
algorithms we will discuss center around the minimization of the
chosen cost function. This depends in turn on our specific
model for describing the data, a typical situation in supervised
learning. Automatizing the search for the minima of the cost function is a
central ingredient in all algorithms. Typical methods which are
employed are various variants of \textbf{gradient} methods. These will be
discussed in more detail later. Again, you'll be surprised to hear that
many practitioners minimize the above function ''by the eye', popularly dubbed as 
'chi by the eye'. That is, change a parameter and see (visually and numerically) that 
the  $\chi^2$ function becomes smaller. 

There are many ways to define the cost function. A simpler approach is to look at the relative difference between the training data and the predicted data, that is we define 
the relative error (why would we prefer the MSE instead of the relative error?) as

\[
\epsilon_{\mathrm{relative}}= \frac{\vert \bm{y} -\bm{\tilde{y}}\vert}{\vert \bm{y}\vert}.
\]

The squared cost function results in an arithmetic mean-unbiased
estimator, and the absolute-value cost function results in a
median-unbiased estimator (in the one-dimensional case, and a
geometric median-unbiased estimator for the multi-dimensional
case). The squared cost function has the disadvantage that it has the tendency
to be dominated by outliers.

We can modify easily the above Python code and plot the relative error instead

















\begin{minted}[fontsize=\fontsize{9pt}{9pt},linenos=false,mathescape,baselinestretch=1.0,fontfamily=tt,xleftmargin=2mm]{python}
import numpy as np
import matplotlib.pyplot as plt
from sklearn.linear_model import LinearRegression

x = np.random.rand(100,1)
y = 5*x+0.01*np.random.randn(100,1)
linreg = LinearRegression()
linreg.fit(x,y)
ypredict = linreg.predict(x)

plt.plot(x, np.abs(ypredict-y)/abs(y), "ro")
plt.axis([0,1.0,0.0, 0.5])
plt.xlabel(r'$x$')
plt.ylabel(r'$\epsilon_{\mathrm{relative}}$')
plt.title(r'Relative error')
plt.show()

\end{minted}


Depending on the parameter in front of the normal distribution, we may
have a small or larger relative error. Try to play around with
different training data sets and study (graphically) the value of the
relative error.

As mentioned above, \textbf{Scikit-Learn} has an impressive functionality.
We can for example extract the values of $\alpha$ and $\beta$ and
their error estimates, or the variance and standard deviation and many
other properties from the statistical data analysis. 

Here we show an
example of the functionality of \textbf{Scikit-Learn}.





























\begin{minted}[fontsize=\fontsize{9pt}{9pt},linenos=false,mathescape,baselinestretch=1.0,fontfamily=tt,xleftmargin=2mm]{python}
import numpy as np 
import matplotlib.pyplot as plt 
from sklearn.linear_model import LinearRegression 
from sklearn.metrics import mean_squared_error, r2_score, mean_squared_log_error, mean_absolute_error

x = np.random.rand(100,1)
y = 2.0+ 5*x+0.5*np.random.randn(100,1)
linreg = LinearRegression()
linreg.fit(x,y)
ypredict = linreg.predict(x)
print('The intercept alpha: \n', linreg.intercept_)
print('Coefficient beta : \n', linreg.coef_)
# The mean squared error                               
print("Mean squared error: %.2f" % mean_squared_error(y, ypredict))
# Explained variance score: 1 is perfect prediction                                 
print('Variance score: %.2f' % r2_score(y, ypredict))
# Mean squared log error                                                        
print('Mean squared log error: %.2f' % mean_squared_log_error(y, ypredict) )
# Mean absolute error                                                           
print('Mean absolute error: %.2f' % mean_absolute_error(y, ypredict))
plt.plot(x, ypredict, "r-")
plt.plot(x, y ,'ro')
plt.axis([0.0,1.0,1.5, 7.0])
plt.xlabel(r'$x$')
plt.ylabel(r'$y$')
plt.title(r'Linear Regression fit ')
plt.show()


\end{minted}

The function \textbf{coef} gives us the parameter $\beta$ of our fit while \textbf{intercept} yields 
$\alpha$. Depending on the constant in front of the normal distribution, we get values near or far from $\alpha =2$ and $\beta =5$. Try to play around with different parameters in front of the normal distribution. The function \textbf{meansquarederror} gives us the mean square error, a risk metric corresponding to the expected value of the squared (quadratic) error or loss defined as
\[ MSE(\bm{y},\bm{\tilde{y}}) = \frac{1}{n}
\sum_{i=0}^{n-1}(y_i-\tilde{y}_i)^2, 
\] 

The smaller the value, the better the fit. Ideally we would like to
have an MSE equal zero.  The attentive reader has probably recognized
this function as being similar to the $\chi^2$ function defined above.

The \textbf{r2score} function computes $R^2$, the coefficient of
determination. It provides a measure of how well future samples are
likely to be predicted by the model. Best possible score is 1.0 and it
can be negative (because the model can be arbitrarily worse). A
constant model that always predicts the expected value of $\bm{y}$,
disregarding the input features, would get a $R^2$ score of $0.0$.

If $\tilde{\bm{y}}_i$ is the predicted value of the $i-th$ sample and $y_i$ is the corresponding true value, then the score $R^2$ is defined as
\[
R^2(\bm{y}, \tilde{\bm{y}}) = 1 - \frac{\sum_{i=0}^{n - 1} (y_i - \tilde{y}_i)^2}{\sum_{i=0}^{n - 1} (y_i - \bar{y})^2},
\]
where we have defined the mean value  of $\bm{y}$ as
\[
\bar{y} =  \frac{1}{n} \sum_{i=0}^{n - 1} y_i.
\]
Another quantity taht we will meet again in our discussions of regression analysis is 
 the mean absolute error (MAE), a risk metric corresponding to the expected value of the absolute error loss or what we call the $l1$-norm loss. In our discussion above we presented the relative error.
The MAE is defined as follows
\[
\text{MAE}(\bm{y}, \bm{\tilde{y}}) = \frac{1}{n} \sum_{i=0}^{n-1} \left| y_i - \tilde{y}_i \right|.
\]
We present the 
squared logarithmic (quadratic) error
\[
\text{MSLE}(\bm{y}, \bm{\tilde{y}}) = \frac{1}{n} \sum_{i=0}^{n - 1} (\log_e (1 + y_i) - \log_e (1 + \tilde{y}_i) )^2,
\]

where $\log_e (x)$ stands for the natural logarithm of $x$. This error
estimate is best to use when targets having exponential growth, such
as population counts, average sales of a commodity over a span of
years etc. 

Finally, another cost function is the Huber cost function used in robust regression.

The rationale behind this possible cost function is its reduced
sensitivity to outliers in the data set. In our discussions on
dimensionality reduction and normalization of data we will meet other
ways of dealing with outliers.

The Huber cost function is defined as
\[
H_{\delta}(\bm{a})=\left\{\begin{array}{cc}\frac{1}{2} \bm{a}^{2}& \text{for }|\bm{a}|\leq \delta\\ \delta (|\bm{a}|-\frac{1}{2}\delta ),&\text{otherwise}.\end{array}\right.
\]
Here $\bm{a}=\bm{y} - \bm{\tilde{y}}$.

We will discuss in more detail these and other functions in the
various lectures.  We conclude this part with another example. Instead
of a linear $x$-dependence we study now a cubic polynomial and use the
polynomial regression analysis tools of scikit-learn.
































\begin{minted}[fontsize=\fontsize{9pt}{9pt},linenos=false,mathescape,baselinestretch=1.0,fontfamily=tt,xleftmargin=2mm]{python}
import matplotlib.pyplot as plt
import numpy as np
import random
from sklearn.linear_model import Ridge
from sklearn.preprocessing import PolynomialFeatures
from sklearn.pipeline import make_pipeline
from sklearn.linear_model import LinearRegression

x=np.linspace(0.02,0.98,200)
noise = np.asarray(random.sample((range(200)),200))
y=x**3*noise
yn=x**3*100
poly3 = PolynomialFeatures(degree=3)
X = poly3.fit_transform(x[:,np.newaxis])
clf3 = LinearRegression()
clf3.fit(X,y)

Xplot=poly3.fit_transform(x[:,np.newaxis])
poly3_plot=plt.plot(x, clf3.predict(Xplot), label='Cubic Fit')
plt.plot(x,yn, color='red', label="True Cubic")
plt.scatter(x, y, label='Data', color='orange', s=15)
plt.legend()
plt.show()

def error(a):
    for i in y:
        err=(y-yn)/yn
    return abs(np.sum(err))/len(err)

print (error(y))

\end{minted}


\noindent\textbf{To our real data: nuclear binding energies. Brief reminder on masses and binding energies.}
Let us now dive into  nuclear physics and remind ourselves briefly about some basic features about binding
energies.  A basic quantity which can be measured for the ground
states of nuclei is the atomic mass $M(N, Z)$ of the neutral atom with
atomic mass number $A$ and charge $Z$. The number of neutrons is $N$. There are indeed several sophisticated experiments worldwide which allow us to measure this quantity to high precision (parts per million even). 

Atomic masses are usually tabulated in terms of the mass excess defined by
\[
\Delta M(N, Z) =  M(N, Z) - uA,
\]
where $u$ is the Atomic Mass Unit 
\[
u = M(^{12}\mathrm{C})/12 = 931.4940954(57) \hspace{0.1cm} \mathrm{MeV}/c^2.
\]
The nucleon masses are
\[
m_p =  1.00727646693(9)u,
\] 
and
\[
m_n = 939.56536(8)\hspace{0.1cm} \mathrm{MeV}/c^2 = 1.0086649156(6)u.
\]

In the \href{{http://nuclearmasses.org/resources_folder/Wang_2017_Chinese_Phys_C_41_030003.pdf}}{2016 mass evaluation of by W.J.Huang, G.Audi, M.Wang, F.G.Kondev, S.Naimi and X.Xu}
there are data on masses and decays of 3437 nuclei.

The nuclear binding energy is defined as the energy required to break
up a given nucleus into its constituent parts of $N$ neutrons and $Z$
protons. In terms of the atomic masses $M(N, Z)$ the binding energy is
defined by

\[
BE(N, Z) = ZM_H c^2 + Nm_n c^2 - M(N, Z)c^2 ,
\]
where $M_H$ is the mass of the hydrogen atom and $m_n$ is the mass of the neutron.
In terms of the mass excess the binding energy is given by
\[
BE(N, Z) = Z\Delta_H c^2 + N\Delta_n c^2 -\Delta(N, Z)c^2 ,
\]
where $\Delta_H c^2 = 7.2890$ MeV and $\Delta_n c^2 = 8.0713$ MeV.

A popular and physically intuitive model which can be used to parametrize 
the experimental binding energies as function of $A$, is the so-called 
\textbf{liquid drop model}. The ansatz is based on the following expression

\[ 
BE(N,Z) = a_1A-a_2A^{2/3}-a_3\frac{Z^2}{A^{1/3}}-a_4\frac{(N-Z)^2}{A},
\]

where $A$ stands for the number of nucleons and the $a_i$s are parameters which are determined by a fit 
to the experimental data.  

To arrive at the above expression we have assumed that we can make the following assumptions:

\begin{itemize}
 \item There is a volume term $a_1A$ proportional with the number of nucleons (the energy is also an extensive quantity). When an assembly of nucleons of the same size is packed together into the smallest volume, each interior nucleon has a certain number of other nucleons in contact with it. This contribution is proportional to the volume.

 \item There is a surface energy term $a_2A^{2/3}$. The assumption here is that a nucleon at the surface of a nucleus interacts with fewer other nucleons than one in the interior of the nucleus and hence its binding energy is less. This surface energy term takes that into account and is therefore negative and is proportional to the surface area.

 \item There is a Coulomb energy term $a_3\frac{Z^2}{A^{1/3}}$. The electric repulsion between each pair of protons in a nucleus yields less binding. 

 \item There is an asymmetry term $a_4\frac{(N-Z)^2}{A}$. This term is associated with the Pauli exclusion principle and reflects the fact that the proton-neutron interaction is more attractive on the average than the neutron-neutron and proton-proton interactions.
\end{itemize}

\noindent
We could also add a so-called pairing term, which is a correction term that
arises from the tendency of proton pairs and neutron pairs to
occur. An even number of particles is more stable than an odd number. 

\noindent\textbf{Organizing our data.}
Let us start with reading and organizing our data. 
We start with the compilation of masses and binding energies from 2016.
After having downloaded this file to our own computer, we are now ready to read the file and start structuring our data.

We start with preparing folders for storing our calculations and the data file over masses and binding energies. We import also various modules that we will find useful in order to present various Machine Learning methods. Here we focus mainly on the functionality of \textbf{scikit-learn}.


































\begin{minted}[fontsize=\fontsize{9pt}{9pt},linenos=false,mathescape,baselinestretch=1.0,fontfamily=tt,xleftmargin=2mm]{python}
# Common imports
import numpy as np
import pandas as pd
import matplotlib.pyplot as plt
import sklearn.linear_model as skl
from sklearn.model_selection import train_test_split
from sklearn.metrics import mean_squared_error, r2_score, mean_absolute_error
import os

# Where to save the figures and data files
PROJECT_ROOT_DIR = "Results"
FIGURE_ID = "Results/FigureFiles"
DATA_ID = "DataFiles/"

if not os.path.exists(PROJECT_ROOT_DIR):
    os.mkdir(PROJECT_ROOT_DIR)

if not os.path.exists(FIGURE_ID):
    os.makedirs(FIGURE_ID)

if not os.path.exists(DATA_ID):
    os.makedirs(DATA_ID)

def image_path(fig_id):
    return os.path.join(FIGURE_ID, fig_id)

def data_path(dat_id):
    return os.path.join(DATA_ID, dat_id)

def save_fig(fig_id):
    plt.savefig(image_path(fig_id) + ".png", format='png')

infile = open(data_path("MassEval2016.dat"),'r')

\end{minted}


Before we proceed, we define also a function for making our plots. You can obviously avoid this and simply set up various \textbf{matplotlib} commands every time you need them. You may however find it convenient to collect all such commands in one function and simply call this function. 












\begin{minted}[fontsize=\fontsize{9pt}{9pt},linenos=false,mathescape,baselinestretch=1.0,fontfamily=tt,xleftmargin=2mm]{python}
from pylab import plt, mpl
plt.style.use('seaborn')
mpl.rcParams['font.family'] = 'serif'

def MakePlot(x,y, styles, labels, axlabels):
    plt.figure(figsize=(10,6))
    for i in range(len(x)):
        plt.plot(x[i], y[i], styles[i], label = labels[i])
        plt.xlabel(axlabels[0])
        plt.ylabel(axlabels[1])
    plt.legend(loc=0)

\end{minted}


Our next step is to read the data on experimental binding energies and
reorganize them as functions of the mass number $A$, the number of
protons $Z$ and neutrons $N$ using \textbf{pandas}.  Before we do this it is
always useful (unless you have a binary file or other types of compressed
data) to actually open the file and simply take a look at it!

In particular, the program that outputs the final nuclear masses is written in Fortran with a specific format. It means that we need to figure out the format and which columns contain the data we are interested in. Pandas comes with a function that reads formatted output. After having admired the file, we are now ready to start massaging it with \textbf{pandas}. The file begins with some basic format information.











\begin{minted}[fontsize=\fontsize{9pt}{9pt},linenos=false,mathescape,baselinestretch=1.0,fontfamily=tt,xleftmargin=2mm]{python}
"""                                                                                                                         
This is taken from the data file of the mass 2016 evaluation.                                                               
All files are 3436 lines long with 124 character per line.                                                                  
       Headers are 39 lines long.                                                                                           
   col 1     :  Fortran character control: 1 = page feed  0 = line feed                                                     
   format    :  a1,i3,i5,i5,i5,1x,a3,a4,1x,f13.5,f11.5,f11.3,f9.3,1x,a2,f11.3,f9.3,1x,i3,1x,f12.5,f11.5                     
   These formats are reflected in the pandas widths variable below, see the statement                                       
   widths=(1,3,5,5,5,1,3,4,1,13,11,11,9,1,2,11,9,1,3,1,12,11,1),                                                            
   Pandas has also a variable header, with length 39 in this case.                                                          
"""

\end{minted}


The data we are interested in are in columns 2, 3, 4 and 11, giving us
the number of neutrons, protons, mass numbers and binding energies,
respectively. We add also for the sake of completeness the element name. The data are in fixed-width formatted lines and we will
covert them into the \textbf{pandas} DataFrame structure.




















\begin{minted}[fontsize=\fontsize{9pt}{9pt},linenos=false,mathescape,baselinestretch=1.0,fontfamily=tt,xleftmargin=2mm]{python}
# Read the experimental data with Pandas
Masses = pd.read_fwf(infile, usecols=(2,3,4,6,11),
              names=('N', 'Z', 'A', 'Element', 'Ebinding'),
              widths=(1,3,5,5,5,1,3,4,1,13,11,11,9,1,2,11,9,1,3,1,12,11,1),
              header=39,
              index_col=False)

# Extrapolated values are indicated by '#' in place of the decimal place, so
# the Ebinding column won't be numeric. Coerce to float and drop these entries.
Masses['Ebinding'] = pd.to_numeric(Masses['Ebinding'], errors='coerce')
Masses = Masses.dropna()
# Convert from keV to MeV.
Masses['Ebinding'] /= 1000

# Group the DataFrame by nucleon number, A.
Masses = Masses.groupby('A')
# Find the rows of the grouped DataFrame with the maximum binding energy.
Masses = Masses.apply(lambda t: t[t.Ebinding==t.Ebinding.max()])

\end{minted}


We have now read in the data, grouped them according to the variables we are interested in. 
We see how easy it is to reorganize the data using \textbf{pandas}. If we
were to do these operations in C/C++ or Fortran, we would have had to
write various functions/subroutines which perform the above
reorganizations for us.  Having reorganized the data, we can now start
to make some simple fits using both the functionalities in \textbf{numpy} and
\textbf{Scikit-Learn} afterwards. 

Now we define five variables which contain
the number of nucleons $A$, the number of protons $Z$ and the number of neutrons $N$, the element name and finally the energies themselves.







\begin{minted}[fontsize=\fontsize{9pt}{9pt},linenos=false,mathescape,baselinestretch=1.0,fontfamily=tt,xleftmargin=2mm]{python}
A = Masses['A']
Z = Masses['Z']
N = Masses['N']
Element = Masses['Element']
Energies = Masses['Ebinding']
print(Masses)

\end{minted}

The next step, and we will define this mathematically later, is to set up the so-called \textbf{design matrix}. We will throughout call this matrix $\bm{X}$.
It has dimensionality $p\times n$, where $n$ is the number of data points and $p$ are the so-called predictors. In our case here they are given by the number of polynomials in $A$ we wish to include in the fit. 








\begin{minted}[fontsize=\fontsize{9pt}{9pt},linenos=false,mathescape,baselinestretch=1.0,fontfamily=tt,xleftmargin=2mm]{python}
# Now we set up the design matrix X
X = np.zeros((len(A),5))
X[:,0] = 1
X[:,1] = A
X[:,2] = A**(2.0/3.0)
X[:,3] = A**(-1.0/3.0)
X[:,4] = A**(-1.0)

\end{minted}

With \textbf{scikitlearn} we are now ready to use linear regression and fit our data.



\begin{minted}[fontsize=\fontsize{9pt}{9pt},linenos=false,mathescape,baselinestretch=1.0,fontfamily=tt,xleftmargin=2mm]{python}
clf = skl.LinearRegression().fit(X, Energies)
fity = clf.predict(X)

\end{minted}

Pretty simple!  
Now we can print measures of how our fit is doing, the coefficients from the fits and plot the final fit together with our data.





















\begin{minted}[fontsize=\fontsize{9pt}{9pt},linenos=false,mathescape,baselinestretch=1.0,fontfamily=tt,xleftmargin=2mm]{python}
# The mean squared error                               
print("Mean squared error: %.2f" % mean_squared_error(Energies, fity))
# Explained variance score: 1 is perfect prediction                                 
print('Variance score: %.2f' % r2_score(Energies, fity))
# Mean absolute error                                                           
print('Mean absolute error: %.2f' % mean_absolute_error(Energies, fity))
print(clf.coef_, clf.intercept_)

Masses['Eapprox']  = fity
# Generate a plot comparing the experimental with the fitted values values.
fig, ax = plt.subplots()
ax.set_xlabel(r'$A = N + Z$')
ax.set_ylabel(r'$E_\mathrm{bind}\,/\mathrm{MeV}$')
ax.plot(Masses['A'], Masses['Ebinding'], alpha=0.7, lw=2,
            label='Ame2016')
ax.plot(Masses['A'], Masses['Eapprox'], alpha=0.7, lw=2, c='m',
            label='Fit')
ax.legend()
save_fig("Masses2016")
plt.show()

\end{minted}


\noindent\textbf{Seeing the wood for the trees.}
As a teaser, let us now see how we can do this with decision trees using \textbf{scikit-learn}. Later we will switch to so-called \textbf{random forests}!
































\begin{minted}[fontsize=\fontsize{9pt}{9pt},linenos=false,mathescape,baselinestretch=1.0,fontfamily=tt,xleftmargin=2mm]{python}

#Decision Tree Regression
from sklearn.tree import DecisionTreeRegressor
regr_1=DecisionTreeRegressor(max_depth=5)
regr_2=DecisionTreeRegressor(max_depth=7)
regr_3=DecisionTreeRegressor(max_depth=9)
regr_1.fit(X, Energies)
regr_2.fit(X, Energies)
regr_3.fit(X, Energies)


y_1 = regr_1.predict(X)
y_2 = regr_2.predict(X)
y_3=regr_3.predict(X)
Masses['Eapprox'] = y_3
# Plot the results
plt.figure()
plt.plot(A, Energies, color="blue", label="Data", linewidth=2)
plt.plot(A, y_1, color="red", label="max_depth=5", linewidth=2)
plt.plot(A, y_2, color="green", label="max_depth=7", linewidth=2)
plt.plot(A, y_3, color="m", label="max_depth=9", linewidth=2)

plt.xlabel("$A$")
plt.ylabel("$E$[MeV]")
plt.title("Decision Tree Regression")
plt.legend()
save_fig("Masses2016Trees")
plt.show()
print(Masses)
print(np.mean( (Energies-y_1)**2))

\end{minted}


\noindent\textbf{And what about using neural networks?}
The \textbf{seaborn} package allows us to visualize data in an efficient way. Note that we use \textbf{scikit-learn}'s multi-layer perceptron (or feed forward neural network) 
functionality.

































\begin{minted}[fontsize=\fontsize{9pt}{9pt},linenos=false,mathescape,baselinestretch=1.0,fontfamily=tt,xleftmargin=2mm]{python}
from sklearn.neural_network import MLPRegressor
from sklearn.metrics import accuracy_score
import seaborn as sns

X_train = X
Y_train = Energies
n_hidden_neurons = 100
epochs = 100
# store models for later use
eta_vals = np.logspace(-5, 1, 7)
lmbd_vals = np.logspace(-5, 1, 7)
# store the models for later use
DNN_scikit = np.zeros((len(eta_vals), len(lmbd_vals)), dtype=object)
train_accuracy = np.zeros((len(eta_vals), len(lmbd_vals)))
sns.set()
for i, eta in enumerate(eta_vals):
    for j, lmbd in enumerate(lmbd_vals):
        dnn = MLPRegressor(hidden_layer_sizes=(n_hidden_neurons), activation='logistic',
                            alpha=lmbd, learning_rate_init=eta, max_iter=epochs)
        dnn.fit(X_train, Y_train)
        DNN_scikit[i][j] = dnn
        train_accuracy[i][j] = dnn.score(X_train, Y_train)

fig, ax = plt.subplots(figsize = (10, 10))
sns.heatmap(train_accuracy, annot=True, ax=ax, cmap="viridis")
ax.set_title("Training Accuracy")
ax.set_ylabel("$\eta$")
ax.set_xlabel("$\lambda$")
plt.show()




\end{minted}




The aim behind these introductory words was to present to you various
Python libraries and their functionalities, in particular libraries like
\textbf{numpy}, \textbf{pandas}, \textbf{xarray} and \textbf{matplotlib} and other that make our life much easier
in handling various data sets and visualizing data. 

Furthermore,
\textbf{Scikit-Learn} allows us with few lines of code to implement popular
Machine Learning algorithms for supervised learning. Later we will meet \textbf{Tensorflow}, a powerful library for deep learning. 
Now it is time to dive more into the details of various methods. We will start with linear regression and try to take a deeper look at what it entails.
\end{frame}

\begin{frame}[plain,fragile]
\frametitle{Why Linear Regression (aka Ordinary Least Squares and family)}

Fitting a continuous function with linear parameterization in terms of the parameters  $\bm{\beta}$.
\begin{itemize}
\item Method of choice for fitting a continuous function!

\item Gives an excellent introduction to central Machine Learning features with \textbf{understandable pedagogical} links to other methods like \textbf{Neural Networks}, \textbf{Support Vector Machines} etc

\item Analytical expression for the fitting parameters $\bm{\beta}$

\item Analytical expressions for statistical propertiers like mean values, variances, confidence intervals and more

\item Analytical relation with probabilistic interpretations 

\item Easy to introduce basic concepts like bias-variance tradeoff, cross-validation, resampling and regularization techniques and many other ML topics

\item Easy to code! And links well with classification problems and logistic regression and neural networks

\item Allows for \textbf{easy} hands-on understanding of gradient descent methods

\item and many more features
\end{itemize}

\noindent
For more discussions of Ridge and Lasso regression, \href{{https://arxiv.org/abs/1509.09169}}{Wessel van Wieringen's} article is highly recommended.
Similarly, \href{{https://arxiv.org/abs/1803.08823}}{Mehta et al's article} is also recommended.
\end{frame}

\begin{frame}[plain,fragile]
\frametitle{Regression analysis, overarching aims}

\begin{block}{}

Regression modeling deals with the description of  the sampling distribution of a given random variable $y$ and how it varies as function of another variable or a set of such variables $\bm{x} =[x_0, x_1,\dots, x_{n-1}]^T$. 
The first variable is called the \textbf{dependent}, the \textbf{outcome} or the \textbf{response} variable while the set of variables $\bm{x}$ is called the independent variable, or the predictor variable or the explanatory variable. 

A regression model aims at finding a likelihood function $p(\bm{y}\vert \bm{x})$, that is the conditional distribution for $\bm{y}$ with a given $\bm{x}$. The estimation of  $p(\bm{y}\vert \bm{x})$ is made using a data set with 
\begin{itemize}
\item $n$ cases $i = 0, 1, 2, \dots, n-1$ 

\item Response (target, dependent or outcome) variable $y_i$ with $i = 0, 1, 2, \dots, n-1$ 

\item $p$ so-called explanatory (independent or predictor) variables $\bm{x}_i=[x_{i0}, x_{i1}, \dots, x_{ip-1}]$ with $i = 0, 1, 2, \dots, n-1$ and explanatory variables running from $0$ to $p-1$. See below for more explicit examples.   
\end{itemize}

\noindent
 The goal of the regression analysis is to extract/exploit relationship between $\bm{y}$ and $\bm{x}$ in or to infer causal dependencies, approximations to the likelihood functions, functional relationships and to make predictions, making fits and many other things.
\end{block}
\end{frame}

\begin{frame}[plain,fragile]
\frametitle{Regression analysis, overarching aims II}

\begin{block}{}

Consider an experiment in which $p$ characteristics of $n$ samples are
measured. The data from this experiment, for various explanatory variables $p$ are normally represented by a matrix  
$\mathbf{X}$.

The matrix $\mathbf{X}$ is called the \emph{design
matrix}. Additional information of the samples is available in the
form of $\bm{y}$ (also as above). The variable $\bm{y}$ is
generally referred to as the \emph{response variable}. The aim of
regression analysis is to explain $\bm{y}$ in terms of
$\bm{X}$ through a functional relationship like $y_i =
f(\mathbf{X}_{i,\ast})$. When no prior knowledge on the form of
$f(\cdot)$ is available, it is common to assume a linear relationship
between $\bm{X}$ and $\bm{y}$. This assumption gives rise to
the \emph{linear regression model} where $\bm{\beta} = [\beta_0, \ldots,
\beta_{p-1}]^{T}$ are the \emph{regression parameters}. 

Linear regression gives us a set of analytical equations for the parameters $\beta_j$.

\end{block}
\end{frame}

\begin{frame}[plain,fragile]
\frametitle{Examples}

\begin{block}{}
In order to understand the relation among the predictors $p$, the set of data $n$ and the target (outcome, output etc) $\bm{y}$,
consider the model we discussed for describing nuclear binding energies. 

There we assumed that we could parametrize the data using a polynomial approximation based on the liquid drop model.
Assuming 
\[
BE(A) = a_0+a_1A+a_2A^{2/3}+a_3A^{-1/3}+a_4A^{-1},
\]
we have five predictors, that is the intercept, the $A$ dependent term, the $A^{2/3}$ term and the $A^{-1/3}$ and $A^{-1}$ terms.
This gives $p=0,1,2,3,4$. Furthermore we have $n$ entries for each predictor. It means that our design matrix is a 
$p\times n$ matrix $\bm{X}$.

Here the predictors are based on a model we have made. A popular data set which is widely encountered in ML applications is the
so-called \href{{https://www.sciencedirect.com/science/article/pii/S0957417407006719?via%3Dihub}}{credit card default data from Taiwan}. The data set contains data on $n=30000$ credit card holders with predictors like gender, marital status, age, profession, education, etc. In total there are $24$ such predictors or attributes leading to a design matrix of dimensionality $24 \times 30000$. This is however a classification problem and we will come back to it when we discuss Logistic Regression. 

\end{block}
\end{frame}

\begin{frame}[plain,fragile]
\frametitle{General linear models}

\begin{block}{}
Before we proceed let us study a case from linear algebra where we aim at fitting a set of data $\bm{y}=[y_0,y_1,\dots,y_{n-1}]$. We could think of these data as a result of an experiment or a complicated numerical experiment. These data are functions of a series of variables $\bm{x}=[x_0,x_1,\dots,x_{n-1}]$, that is $y_i = y(x_i)$ with $i=0,1,2,\dots,n-1$. The variables $x_i$ could represent physical quantities like time, temperature, position etc. We assume that $y(x)$ is a smooth function. 

Since obtaining these data points may not be trivial, we want to use these data to fit a function which can allow us to make predictions for values of $y$ which are not in the present set. The perhaps simplest approach is to assume we can parametrize our function in terms of a polynomial of degree $n-1$ with $n$ points, that is
\[
y=y(x) \rightarrow y(x_i)=\tilde{y}_i+\epsilon_i=\sum_{j=0}^{n-1} \beta_j x_i^j+\epsilon_i,
\]
where $\epsilon_i$ is the error in our approximation. 

\end{block}
\end{frame}

\begin{frame}[plain,fragile]
\frametitle{Rewriting the fitting procedure as a linear algebra problem}

\begin{block}{}
For every set of values $y_i,x_i$ we have thus the corresponding set of equations
\begin{align*}
y_0&=\beta_0+\beta_1x_0^1+\beta_2x_0^2+\dots+\beta_{n-1}x_0^{n-1}+\epsilon_0\\
y_1&=\beta_0+\beta_1x_1^1+\beta_2x_1^2+\dots+\beta_{n-1}x_1^{n-1}+\epsilon_1\\
y_2&=\beta_0+\beta_1x_2^1+\beta_2x_2^2+\dots+\beta_{n-1}x_2^{n-1}+\epsilon_2\\
\dots & \dots \\
y_{n-1}&=\beta_0+\beta_1x_{n-1}^1+\beta_2x_{n-1}^2+\dots+\beta_{n-1}x_{n-1}^{n-1}+\epsilon_{n-1}.\\
\end{align*}
\end{block}
\end{frame}

\begin{frame}[plain,fragile]
\frametitle{Rewriting the fitting procedure as a linear algebra problem, more details}

\begin{block}{}
Defining the vectors
\[
\bm{y} = [y_0,y_1, y_2,\dots, y_{n-1}]^T,
\]
and
\[
\bm{\beta} = [\beta_0,\beta_1, \beta_2,\dots, \beta_{n-1}]^T,
\]
and
\[
\bm{\epsilon} = [\epsilon_0,\epsilon_1, \epsilon_2,\dots, \epsilon_{n-1}]^T,
\]
and the design matrix
\[
\bm{X}=
\begin{bmatrix} 
1& x_{0}^1 &x_{0}^2& \dots & \dots &x_{0}^{n-1}\\
1& x_{1}^1 &x_{1}^2& \dots & \dots &x_{1}^{n-1}\\
1& x_{2}^1 &x_{2}^2& \dots & \dots &x_{2}^{n-1}\\                      
\dots& \dots &\dots& \dots & \dots &\dots\\
1& x_{n-1}^1 &x_{n-1}^2& \dots & \dots &x_{n-1}^{n-1}\\
\end{bmatrix} 
\]
we can rewrite our equations as
\[
\bm{y} = \bm{X}\bm{\beta}+\bm{\epsilon}.
\]
The above design matrix is called a \href{{https://en.wikipedia.org/wiki/Vandermonde_matrix}}{Vandermonde matrix}.
\end{block}
\end{frame}

\begin{frame}[plain,fragile]
\frametitle{Generalizing the fitting procedure as a linear algebra problem}

\begin{block}{}

We are obviously not limited to the above polynomial expansions.  We
could replace the various powers of $x$ with elements of Fourier
series or instead of $x_i^j$ we could have $\cos{(j x_i)}$ or $\sin{(j
x_i)}$, or time series or other orthogonal functions.  For every set
of values $y_i,x_i$ we can then generalize the equations to

\begin{align*}
y_0&=\beta_0x_{00}+\beta_1x_{01}+\beta_2x_{02}+\dots+\beta_{n-1}x_{0n-1}+\epsilon_0\\
y_1&=\beta_0x_{10}+\beta_1x_{11}+\beta_2x_{12}+\dots+\beta_{n-1}x_{1n-1}+\epsilon_1\\
y_2&=\beta_0x_{20}+\beta_1x_{21}+\beta_2x_{22}+\dots+\beta_{n-1}x_{2n-1}+\epsilon_2\\
\dots & \dots \\
y_{i}&=\beta_0x_{i0}+\beta_1x_{i1}+\beta_2x_{i2}+\dots+\beta_{n-1}x_{in-1}+\epsilon_i\\
\dots & \dots \\
y_{n-1}&=\beta_0x_{n-1,0}+\beta_1x_{n-1,2}+\beta_2x_{n-1,2}+\dots+\beta_{n-1}x_{n-1,n-1}+\epsilon_{n-1}.\\
\end{align*}

\textbf{Note that we have $p=n$ here. The matrix is symmetric. This is generally not the case!}
\end{block}
\end{frame}

\begin{frame}[plain,fragile]
\frametitle{Generalizing the fitting procedure as a linear algebra problem}

\begin{block}{}
We redefine in turn the matrix $\bm{X}$ as
\[
\bm{X}=
\begin{bmatrix} 
x_{00}& x_{01} &x_{02}& \dots & \dots &x_{0,n-1}\\
x_{10}& x_{11} &x_{12}& \dots & \dots &x_{1,n-1}\\
x_{20}& x_{21} &x_{22}& \dots & \dots &x_{2,n-1}\\                      
\dots& \dots &\dots& \dots & \dots &\dots\\
x_{n-1,0}& x_{n-1,1} &x_{n-1,2}& \dots & \dots &x_{n-1,n-1}\\
\end{bmatrix} 
\]
and without loss of generality we rewrite again  our equations as
\[
\bm{y} = \bm{X}\bm{\beta}+\bm{\epsilon}.
\]
The left-hand side of this equation is kwown. Our error vector $\bm{\epsilon}$ and the parameter vector $\bm{\beta}$ are our unknow quantities. How can we obtain the optimal set of $\beta_i$ values? 
\end{block}
\end{frame}

\begin{frame}[plain,fragile]
\frametitle{Optimizing our parameters}

\begin{block}{}
We have defined the matrix $\bm{X}$ via the equations
\begin{align*}
y_0&=\beta_0x_{00}+\beta_1x_{01}+\beta_2x_{02}+\dots+\beta_{n-1}x_{0n-1}+\epsilon_0\\
y_1&=\beta_0x_{10}+\beta_1x_{11}+\beta_2x_{12}+\dots+\beta_{n-1}x_{1n-1}+\epsilon_1\\
y_2&=\beta_0x_{20}+\beta_1x_{21}+\beta_2x_{22}+\dots+\beta_{n-1}x_{2n-1}+\epsilon_1\\
\dots & \dots \\
y_{i}&=\beta_0x_{i0}+\beta_1x_{i1}+\beta_2x_{i2}+\dots+\beta_{n-1}x_{in-1}+\epsilon_1\\
\dots & \dots \\
y_{n-1}&=\beta_0x_{n-1,0}+\beta_1x_{n-1,2}+\beta_2x_{n-1,2}+\dots+\beta_{n-1}x_{n-1,n-1}+\epsilon_{n-1}.\\
\end{align*}

As we noted above, we stayed with a system with the design matrix 
 $\bm{X}\in {\mathbb{R}}^{n\times n}$, that is we have $p=n$. For reasons to come later (algorithmic arguments) we will hereafter define 
our matrix as $\bm{X}\in {\mathbb{R}}^{n\times p}$, with the predictors refering to the column numbers and the entries $n$ being the row elements.

\end{block}
\end{frame}

\begin{frame}[plain,fragile]
\frametitle{Our model for the nuclear binding energies}

In our \href{{https://compphysics.github.io/MachineLearning/doc/pub/How2ReadData/html/How2ReadData.html}}{introductory notes} we looked at the so-called \href{{https://en.wikipedia.org/wiki/Semi-empirical_mass_formula}}{liquid drop model}. Let us remind ourselves about what we did by looking at the code.

We restate the parts of the code we are most interested in.






































































\begin{minted}[fontsize=\fontsize{9pt}{9pt},linenos=false,mathescape,baselinestretch=1.0,fontfamily=tt,xleftmargin=2mm]{python}
# Common imports
import numpy as np
import pandas as pd
import matplotlib.pyplot as plt
from IPython.display import display
import os

# Where to save the figures and data files
PROJECT_ROOT_DIR = "Results"
FIGURE_ID = "Results/FigureFiles"
DATA_ID = "DataFiles/"

if not os.path.exists(PROJECT_ROOT_DIR):
    os.mkdir(PROJECT_ROOT_DIR)

if not os.path.exists(FIGURE_ID):
    os.makedirs(FIGURE_ID)

if not os.path.exists(DATA_ID):
    os.makedirs(DATA_ID)

def image_path(fig_id):
    return os.path.join(FIGURE_ID, fig_id)

def data_path(dat_id):
    return os.path.join(DATA_ID, dat_id)

def save_fig(fig_id):
    plt.savefig(image_path(fig_id) + ".png", format='png')

infile = open(data_path("MassEval2016.dat"),'r')


# Read the experimental data with Pandas
Masses = pd.read_fwf(infile, usecols=(2,3,4,6,11),
              names=('N', 'Z', 'A', 'Element', 'Ebinding'),
              widths=(1,3,5,5,5,1,3,4,1,13,11,11,9,1,2,11,9,1,3,1,12,11,1),
              header=39,
              index_col=False)

# Extrapolated values are indicated by '#' in place of the decimal place, so
# the Ebinding column won't be numeric. Coerce to float and drop these entries.
Masses['Ebinding'] = pd.to_numeric(Masses['Ebinding'], errors='coerce')
Masses = Masses.dropna()
# Convert from keV to MeV.
Masses['Ebinding'] /= 1000

# Group the DataFrame by nucleon number, A.
Masses = Masses.groupby('A')
# Find the rows of the grouped DataFrame with the maximum binding energy.
Masses = Masses.apply(lambda t: t[t.Ebinding==t.Ebinding.max()])
A = Masses['A']
Z = Masses['Z']
N = Masses['N']
Element = Masses['Element']
Energies = Masses['Ebinding']

# Now we set up the design matrix X
X = np.zeros((len(A),5))
X[:,0] = 1
X[:,1] = A
X[:,2] = A**(2.0/3.0)
X[:,3] = A**(-1.0/3.0)
X[:,4] = A**(-1.0)
# Then nice printout using pandas
DesignMatrix = pd.DataFrame(X)
DesignMatrix.index = A
DesignMatrix.columns = ['1', 'A', 'A^(2/3)', 'A^(-1/3)', '1/A']
display(DesignMatrix)

\end{minted}


With $\bm{\beta}\in {\mathbb{R}}^{p\times 1}$, it means that we will hereafter write our equations for the approximation as
\[
\bm{\tilde{y}}= \bm{X}\bm{\beta},
\]
throughout these lectures.
\end{frame}

\begin{frame}[plain,fragile]
\frametitle{Optimizing our parameters, more details}

\begin{block}{}
With the above we use the design matrix to define the approximation $\bm{\tilde{y}}$ via the unknown quantity $\bm{\beta}$ as
\[
\bm{\tilde{y}}= \bm{X}\bm{\beta},
\]
and in order to find the optimal parameters $\beta_i$ instead of solving the above linear algebra problem, we define a function which gives a measure of the spread between the values $y_i$ (which represent hopefully the exact values) and the parameterized values $\tilde{y}_i$, namely
\[
C(\bm{\beta})=\frac{1}{n}\sum_{i=0}^{n-1}\left(y_i-\tilde{y}_i\right)^2=\frac{1}{n}\left\{\left(\bm{y}-\bm{\tilde{y}}\right)^T\left(\bm{y}-\bm{\tilde{y}}\right)\right\},
\]
or using the matrix $\bm{X}$ and in a more compact matrix-vector notation as
\[
C(\bm{\beta})=\frac{1}{n}\left\{\left(\bm{y}-\bm{X}\bm{\beta}\right)^T\left(\bm{y}-\bm{X}\bm{\beta}\right)\right\}.
\]
This function is one possible way to define the so-called cost function.

It is also common to define
the function $C$ as

\[
C(\bm{\beta})=\frac{1}{2n}\sum_{i=0}^{n-1}\left(y_i-\tilde{y}_i\right)^2,
\]
since when taking the first derivative with respect to the unknown parameters $\beta$, the factor of $2$ cancels out. 
\end{block}
\end{frame}

\begin{frame}[plain,fragile]
\frametitle{Interpretations and optimizing our parameters}

\begin{block}{}

The function 
\[
C(\bm{\beta})=\frac{1}{n}\left\{\left(\bm{y}-\bm{X}\bm{\beta}\right)^T\left(\bm{y}-\bm{X}\bm{\beta}\right)\right\},
\]
can be linked to the variance of the quantity $y_i$ if we interpret the latter as the mean value. 
When linking (see the discussion below) with the maximum likelihood approach below, we will indeed interpret $y_i$ as a mean value
\[
y_{i}=\langle y_i \rangle = \beta_0x_{i,0}+\beta_1x_{i,1}+\beta_2x_{i,2}+\dots+\beta_{n-1}x_{i,n-1}+\epsilon_i,
\]

where $\langle y_i \rangle$ is the mean value. Keep in mind also that
till now we have treated $y_i$ as the exact value. Normally, the
response (dependent or outcome) variable $y_i$ the outcome of a
numerical experiment or another type of experiment and is thus only an
approximation to the true value. It is then always accompanied by an
error estimate, often limited to a statistical error estimate given by
the standard deviation discussed earlier. In the discussion here we
will treat $y_i$ as our exact value for the response variable.

In order to find the parameters $\beta_i$ we will then minimize the spread of $C(\bm{\beta})$, that is we are going to solve the problem
\[
{\displaystyle \min_{\bm{\beta}\in
{\mathbb{R}}^{p}}}\frac{1}{n}\left\{\left(\bm{y}-\bm{X}\bm{\beta}\right)^T\left(\bm{y}-\bm{X}\bm{\beta}\right)\right\}.
\]
In practical terms it means we will require
\[
\frac{\partial C(\bm{\beta})}{\partial \beta_j} = \frac{\partial }{\partial \beta_j}\left[ \frac{1}{n}\sum_{i=0}^{n-1}\left(y_i-\beta_0x_{i,0}-\beta_1x_{i,1}-\beta_2x_{i,2}-\dots-\beta_{n-1}x_{i,n-1}\right)^2\right]=0, 
\]
which results in
\[
\frac{\partial C(\bm{\beta})}{\partial \beta_j} = -\frac{2}{n}\left[ \sum_{i=0}^{n-1}x_{ij}\left(y_i-\beta_0x_{i,0}-\beta_1x_{i,1}-\beta_2x_{i,2}-\dots-\beta_{n-1}x_{i,n-1}\right)\right]=0, 
\]
or in a matrix-vector form as
\[
\frac{\partial C(\bm{\beta})}{\partial \bm{\beta}} = 0 = \bm{X}^T\left( \bm{y}-\bm{X}\bm{\beta}\right).  
\]

\end{block}
\end{frame}

\begin{frame}[plain,fragile]
\frametitle{Interpretations and optimizing our parameters}

\begin{block}{}
We can rewrite
\[
\frac{\partial C(\bm{\beta})}{\partial \bm{\beta}} = 0 = \bm{X}^T\left( \bm{y}-\bm{X}\bm{\beta}\right),  
\]
as
\[
\bm{X}^T\bm{y} = \bm{X}^T\bm{X}\bm{\beta},  
\]
and if the matrix $\bm{X}^T\bm{X}$ is invertible we have the solution
\[
\bm{\beta} =\left(\bm{X}^T\bm{X}\right)^{-1}\bm{X}^T\bm{y}.
\]

We note also that since our design matrix is defined as $\bm{X}\in
{\mathbb{R}}^{n\times p}$, the product $\bm{X}^T\bm{X} \in
{\mathbb{R}}^{p\times p}$.  In the above case we have that $p \ll n$,
in our case $p=5$ meaning that we end up with inverting a small
$5\times 5$ matrix. This is a rather common situation, in many cases we end up with low-dimensional
matrices to invert. The methods discussed here and for many other
supervised learning algorithms like classification with logistic
regression or support vector machines, exhibit dimensionalities which
allow for the usage of direct linear algebra methods such as \textbf{LU} decomposition or \textbf{Singular Value Decomposition} (SVD) for finding the inverse of the matrix
$\bm{X}^T\bm{X}$. 
\end{block}

\begin{block}{}
\textbf{Small question}: Do you think the example we have at hand here (the nuclear binding energies) can lead to problems in inverting the matrix  $\bm{X}^T\bm{X}$? What kind of problems can we expect? 
\end{block}
\end{frame}

\begin{frame}[plain,fragile]
\frametitle{Some useful matrix and vector expressions}

The following matrix and vector relation will be useful here and for the rest of the course. Vectors are always written as boldfaced lower case letters and 
matrices as upper case boldfaced letters. 

\[
\frac{\partial (\bm{b}^T\bm{a})}{\partial \bm{a}} = \bm{b},
\]
\[
\frac{\partial (\bm{a}^T\bm{A}\bm{a})}{\partial \bm{a}} = (\bm{A}+\bm{A}^T)\bm{a},
\]
\[
\frac{\partial tr(\bm{B}\bm{A})}{\partial \bm{A}} = \bm{B}^T,
\]
\[
\frac{\partial \log{\vert\bm{A}\vert}}{\partial \bm{A}} = (\bm{A}^{-1})^T.
\]
\end{frame}

\begin{frame}[plain,fragile]
\frametitle{Interpretations and optimizing our parameters}

\begin{block}{}
The residuals $\bm{\epsilon}$ are in turn given by
\[
\bm{\epsilon} = \bm{y}-\bm{\tilde{y}} = \bm{y}-\bm{X}\bm{\beta},
\]
and with 
\[
\bm{X}^T\left( \bm{y}-\bm{X}\bm{\beta}\right)= 0, 
\]
we have
\[
\bm{X}^T\bm{\epsilon}=\bm{X}^T\left( \bm{y}-\bm{X}\bm{\beta}\right)= 0, 
\]
meaning that the solution for $\bm{\beta}$ is the one which minimizes the residuals.  Later we will link this with the maximum likelihood approach.

\end{block}

Let us now return to our nuclear binding energies and simply code the above equations.
\end{frame}

\begin{frame}[plain,fragile]
\frametitle{Own code for Ordinary Least Squares}

It is rather straightforward to implement the matrix inversion and obtain the parameters $\bm{\beta}$. After having defined the matrix $\bm{X}$ we simply need to 
write





\begin{minted}[fontsize=\fontsize{9pt}{9pt},linenos=false,mathescape,baselinestretch=1.0,fontfamily=tt,xleftmargin=2mm]{python}
# matrix inversion to find beta
beta = np.linalg.inv(X.T.dot(X)).dot(X.T).dot(Energies)
# and then make the prediction
ytilde = X @ beta

\end{minted}

Alternatively, you can use the least squares functionality in \textbf{Numpy} as



\begin{minted}[fontsize=\fontsize{9pt}{9pt},linenos=false,mathescape,baselinestretch=1.0,fontfamily=tt,xleftmargin=2mm]{python}
fit = np.linalg.lstsq(X, Energies, rcond =None)[0]
ytildenp = np.dot(fit,X.T)

\end{minted}


And finally we plot our fit with and compare with data













\begin{minted}[fontsize=\fontsize{9pt}{9pt},linenos=false,mathescape,baselinestretch=1.0,fontfamily=tt,xleftmargin=2mm]{python}
Masses['Eapprox']  = ytilde
# Generate a plot comparing the experimental with the fitted values values.
fig, ax = plt.subplots()
ax.set_xlabel(r'$A = N + Z$')
ax.set_ylabel(r'$E_\mathrm{bind}\,/\mathrm{MeV}$')
ax.plot(Masses['A'], Masses['Ebinding'], alpha=0.7, lw=2,
            label='Ame2016')
ax.plot(Masses['A'], Masses['Eapprox'], alpha=0.7, lw=2, c='m',
            label='Fit')
ax.legend()
save_fig("Masses2016OLS")
plt.show()

\end{minted}
\end{frame}

\begin{frame}[plain,fragile]
\frametitle{Adding error analysis and training set up}

We can easily test our fit by computing the $R2$ score that we discussed in connection with the functionality of \textbf{Scikit-Learn} in the introductory slides.
Since we are not using \textbf{Scikit-Learn} here we can define our own $R2$ function as



\begin{minted}[fontsize=\fontsize{9pt}{9pt},linenos=false,mathescape,baselinestretch=1.0,fontfamily=tt,xleftmargin=2mm]{python}
def R2(y_data, y_model):
    return 1 - np.sum((y_data - y_model) ** 2) / np.sum((y_data - np.mean(y_data)) ** 2)

\end{minted}

and we would be using it as 


\begin{minted}[fontsize=\fontsize{9pt}{9pt},linenos=false,mathescape,baselinestretch=1.0,fontfamily=tt,xleftmargin=2mm]{python}
print(R2(Energies,ytilde))

\end{minted}


We can easily add our \textbf{MSE} score as






\begin{minted}[fontsize=\fontsize{9pt}{9pt},linenos=false,mathescape,baselinestretch=1.0,fontfamily=tt,xleftmargin=2mm]{python}
def MSE(y_data,y_model):
    n = np.size(y_model)
    return np.sum((y_data-y_model)**2)/n

print(MSE(Energies,ytilde))

\end{minted}

and finally the relative error as




\begin{minted}[fontsize=\fontsize{9pt}{9pt},linenos=false,mathescape,baselinestretch=1.0,fontfamily=tt,xleftmargin=2mm]{python}
def RelativeError(y_data,y_model):
    return abs((y_data-y_model)/y_data)
print(RelativeError(Energies, ytilde))

\end{minted}
\end{frame}

\begin{frame}[plain,fragile]
\frametitle{The $\chi^2$ function}

\begin{block}{}

Normally, the response (dependent or outcome) variable $y_i$ is the
outcome of a numerical experiment or another type of experiment and is
thus only an approximation to the true value. It is then always
accompanied by an error estimate, often limited to a statistical error
estimate given by the standard deviation discussed earlier. In the
discussion here we will treat $y_i$ as our exact value for the
response variable.

Introducing the standard deviation $\sigma_i$ for each measurement
$y_i$, we define now the $\chi^2$ function (omitting the $1/n$ term)
as

\[
\chi^2(\bm{\beta})=\frac{1}{n}\sum_{i=0}^{n-1}\frac{\left(y_i-\tilde{y}_i\right)^2}{\sigma_i^2}=\frac{1}{n}\left\{\left(\bm{y}-\bm{\tilde{y}}\right)^T\frac{1}{\bm{\Sigma^2}}\left(\bm{y}-\bm{\tilde{y}}\right)\right\},
\]
where the matrix $\bm{\Sigma}$ is a diagonal matrix with $\sigma_i$ as matrix elements. 

\end{block}
\end{frame}

\begin{frame}[plain,fragile]
\frametitle{The $\chi^2$ function}

\begin{block}{}

In order to find the parameters $\beta_i$ we will then minimize the spread of $\chi^2(\bm{\beta})$ by requiring
\[
\frac{\partial \chi^2(\bm{\beta})}{\partial \beta_j} = \frac{\partial }{\partial \beta_j}\left[ \frac{1}{n}\sum_{i=0}^{n-1}\left(\frac{y_i-\beta_0x_{i,0}-\beta_1x_{i,1}-\beta_2x_{i,2}-\dots-\beta_{n-1}x_{i,n-1}}{\sigma_i}\right)^2\right]=0, 
\]
which results in
\[
\frac{\partial \chi^2(\bm{\beta})}{\partial \beta_j} = -\frac{2}{n}\left[ \sum_{i=0}^{n-1}\frac{x_{ij}}{\sigma_i}\left(\frac{y_i-\beta_0x_{i,0}-\beta_1x_{i,1}-\beta_2x_{i,2}-\dots-\beta_{n-1}x_{i,n-1}}{\sigma_i}\right)\right]=0, 
\]
or in a matrix-vector form as
\[
\frac{\partial \chi^2(\bm{\beta})}{\partial \bm{\beta}} = 0 = \bm{A}^T\left( \bm{b}-\bm{A}\bm{\beta}\right).  
\]
where we have defined the matrix $\bm{A} =\bm{X}/\bm{\Sigma}$ with matrix elements $a_{ij} = x_{ij}/\sigma_i$ and the vector $\bm{b}$ with elements $b_i = y_i/\sigma_i$.   
\end{block}
\end{frame}

\begin{frame}[plain,fragile]
\frametitle{The $\chi^2$ function}

\begin{block}{}

We can rewrite
\[
\frac{\partial \chi^2(\bm{\beta})}{\partial \bm{\beta}} = 0 = \bm{A}^T\left( \bm{b}-\bm{A}\bm{\beta}\right),  
\]
as
\[
\bm{A}^T\bm{b} = \bm{A}^T\bm{A}\bm{\beta},  
\]
and if the matrix $\bm{A}^T\bm{A}$ is invertible we have the solution
\[
\bm{\beta} =\left(\bm{A}^T\bm{A}\right)^{-1}\bm{A}^T\bm{b}.
\]
\end{block}
\end{frame}

\begin{frame}[plain,fragile]
\frametitle{The $\chi^2$ function}

\begin{block}{}

If we then introduce the matrix
\[
\bm{H} =  \left(\bm{A}^T\bm{A}\right)^{-1},
\]
we have then the following expression for the parameters $\beta_j$ (the matrix elements of $\bm{H}$ are $h_{ij}$)
\[
\beta_j = \sum_{k=0}^{p-1}h_{jk}\sum_{i=0}^{n-1}\frac{y_i}{\sigma_i}\frac{x_{ik}}{\sigma_i} = \sum_{k=0}^{p-1}h_{jk}\sum_{i=0}^{n-1}b_ia_{ik}
\]
We state without proof the expression for the uncertainty  in the parameters $\beta_j$ as (we leave this as an exercise)
\[
\sigma^2(\beta_j) = \sum_{i=0}^{n-1}\sigma_i^2\left( \frac{\partial \beta_j}{\partial y_i}\right)^2, 
\]
resulting in 
\[
\sigma^2(\beta_j) = \left(\sum_{k=0}^{p-1}h_{jk}\sum_{i=0}^{n-1}a_{ik}\right)\left(\sum_{l=0}^{p-1}h_{jl}\sum_{m=0}^{n-1}a_{ml}\right) = h_{jj}!
\]
\end{block}
\end{frame}

\begin{frame}[plain,fragile]
\frametitle{The $\chi^2$ function}

\begin{block}{}
The first step here is to approximate the function $y$ with a first-order polynomial, that is we write
\[
y=y(x) \rightarrow y(x_i) \approx \beta_0+\beta_1 x_i.
\]
By computing the derivatives of $\chi^2$ with respect to $\beta_0$ and $\beta_1$ show that these are given by
\[
\frac{\partial \chi^2(\bm{\beta})}{\partial \beta_0} = -2\left[ \frac{1}{n}\sum_{i=0}^{n-1}\left(\frac{y_i-\beta_0-\beta_1x_{i}}{\sigma_i^2}\right)\right]=0,
\]
and
\[
\frac{\partial \chi^2(\bm{\beta})}{\partial \beta_1} = -\frac{2}{n}\left[ \sum_{i=0}^{n-1}x_i\left(\frac{y_i-\beta_0-\beta_1x_{i}}{\sigma_i^2}\right)\right]=0.
\]
\end{block}
\end{frame}

\begin{frame}[plain,fragile]
\frametitle{The $\chi^2$ function}

\begin{block}{}

For a linear fit (a first-order polynomial) we don't need to invert a matrix!!  
Defining
\[
\gamma =  \sum_{i=0}^{n-1}\frac{1}{\sigma_i^2},
\]

\[
\gamma_x =  \sum_{i=0}^{n-1}\frac{x_{i}}{\sigma_i^2},
\]

\[
\gamma_y = \sum_{i=0}^{n-1}\left(\frac{y_i}{\sigma_i^2}\right),
\]

\[
\gamma_{xx} =  \sum_{i=0}^{n-1}\frac{x_ix_{i}}{\sigma_i^2},
\]

\[
\gamma_{xy} = \sum_{i=0}^{n-1}\frac{y_ix_{i}}{\sigma_i^2},
\]

we obtain

\[
\beta_0 = \frac{\gamma_{xx}\gamma_y-\gamma_x\gamma_y}{\gamma\gamma_{xx}-\gamma_x^2},
\]

\[
\beta_1 = \frac{\gamma_{xy}\gamma-\gamma_x\gamma_y}{\gamma\gamma_{xx}-\gamma_x^2}.
\]

This approach (different linear and non-linear regression) suffers
often from both being underdetermined and overdetermined in the
unknown coefficients $\beta_i$.  A better approach is to use the
Singular Value Decomposition (SVD) method discussed next week.

\end{block}
\end{frame}

\begin{frame}[plain,fragile]
\frametitle{Fitting an Equation of State for Dense Nuclear Matter}

Before we continue, let us introduce yet another example. We are going to fit the
nuclear equation of state using results from many-body calculations.
The equation of state we have made available here, as function of
density, has been derived using modern nucleon-nucleon potentials with
\href{{https://www.sciencedirect.com/science/article/pii/S0370157399001106}}{the addition of three-body
forces}. This
time the file is presented as a standard \textbf{csv} file.

The beginning of the Python code here is similar to what you have seen
before, with the same initializations and declarations. We use also
\textbf{pandas} again, rather extensively in order to organize our data.

The difference now is that we use \textbf{Scikit-Learn's} regression tools
instead of our own matrix inversion implementation. Furthermore, we
sneak in \textbf{Ridge} regression (to be discussed below) which includes a
hyperparameter $\lambda$, also to be explained below.
\end{frame}

\begin{frame}[plain,fragile]
\frametitle{The code}

\begin{minted}[fontsize=\fontsize{9pt}{9pt},linenos=false,mathescape,baselinestretch=1.0,fontfamily=tt,xleftmargin=2mm]{python}
# Common imports
import os
import numpy as np
import pandas as pd
import matplotlib.pyplot as plt
import matplotlib.pyplot as plt
import sklearn.linear_model as skl
from sklearn.metrics import mean_squared_error, r2_score, mean_absolute_error

# Where to save the figures and data files
PROJECT_ROOT_DIR = "Results"
FIGURE_ID = "Results/FigureFiles"
DATA_ID = "DataFiles/"

if not os.path.exists(PROJECT_ROOT_DIR):
    os.mkdir(PROJECT_ROOT_DIR)

if not os.path.exists(FIGURE_ID):
    os.makedirs(FIGURE_ID)

if not os.path.exists(DATA_ID):
    os.makedirs(DATA_ID)

def image_path(fig_id):
    return os.path.join(FIGURE_ID, fig_id)

def data_path(dat_id):
    return os.path.join(DATA_ID, dat_id)

def save_fig(fig_id):
    plt.savefig(image_path(fig_id) + ".png", format='png')

infile = open(data_path("EoS.csv"),'r')

# Read the EoS data as  csv file and organize the data into two arrays with density and energies
EoS = pd.read_csv(infile, names=('Density', 'Energy'))
EoS['Energy'] = pd.to_numeric(EoS['Energy'], errors='coerce')
EoS = EoS.dropna()
Energies = EoS['Energy']
Density = EoS['Density']
#  The design matrix now as function of various polytrops
X = np.zeros((len(Density),4))
X[:,3] = Density**(4.0/3.0)
X[:,2] = Density
X[:,1] = Density**(2.0/3.0)
X[:,0] = 1

# We use now Scikit-Learn's linear regressor and ridge regressor
# OLS part
clf = skl.LinearRegression().fit(X, Energies)
ytilde = clf.predict(X)
EoS['Eols']  = ytilde
# The mean squared error                               
print("Mean squared error: %.2f" % mean_squared_error(Energies, ytilde))
# Explained variance score: 1 is perfect prediction                                 
print('Variance score: %.2f' % r2_score(Energies, ytilde))
# Mean absolute error                                                           
print('Mean absolute error: %.2f' % mean_absolute_error(Energies, ytilde))
print(clf.coef_, clf.intercept_)

# The Ridge regression with a hyperparameter lambda = 0.1
_lambda = 0.1
clf_ridge = skl.Ridge(alpha=_lambda).fit(X, Energies)
yridge = clf_ridge.predict(X)
EoS['Eridge']  = yridge
# The mean squared error                               
print("Mean squared error: %.2f" % mean_squared_error(Energies, yridge))
# Explained variance score: 1 is perfect prediction                                 
print('Variance score: %.2f' % r2_score(Energies, yridge))
# Mean absolute error                                                           
print('Mean absolute error: %.2f' % mean_absolute_error(Energies, yridge))
print(clf_ridge.coef_, clf_ridge.intercept_)

fig, ax = plt.subplots()
ax.set_xlabel(r'$\rho[\mathrm{fm}^{-3}]$')
ax.set_ylabel(r'Energy per particle')
ax.plot(EoS['Density'], EoS['Energy'], alpha=0.7, lw=2,
            label='Theoretical data')
ax.plot(EoS['Density'], EoS['Eols'], alpha=0.7, lw=2, c='m',
            label='OLS')
ax.plot(EoS['Density'], EoS['Eridge'], alpha=0.7, lw=2, c='g',
            label='Ridge $\lambda = 0.1$')
ax.legend()
save_fig("EoSfitting")
plt.show()

\end{minted}


The above simple polynomial in density $\rho$ gives an excellent fit
to the data. 

We note also that there is a small deviation between the
standard OLS and the Ridge regression at higher densities. We discuss this in more detail
below.
\end{frame}

\begin{frame}[plain,fragile]
\frametitle{Splitting our Data in Training and Test data}

It is normal in essentially all Machine Learning studies to split the
data in a training set and a test set (sometimes also an additional
validation set).  \textbf{Scikit-Learn} has an own function for this. There
is no explicit recipe for how much data should be included as training
data and say test data.  An accepted rule of thumb is to use
approximately $2/3$ to $4/5$ of the data as training data. We will
postpone a discussion of this splitting to the end of these notes and
our discussion of the so-called \textbf{bias-variance} tradeoff. Here we
limit ourselves to repeat the above equation of state fitting example
but now splitting the data into a training set and a test set.


































































\begin{minted}[fontsize=\fontsize{9pt}{9pt},linenos=false,mathescape,baselinestretch=1.0,fontfamily=tt,xleftmargin=2mm]{python}
import os
import numpy as np
import pandas as pd
import matplotlib.pyplot as plt
from sklearn.model_selection import train_test_split
# Where to save the figures and data files
PROJECT_ROOT_DIR = "Results"
FIGURE_ID = "Results/FigureFiles"
DATA_ID = "DataFiles/"

if not os.path.exists(PROJECT_ROOT_DIR):
    os.mkdir(PROJECT_ROOT_DIR)

if not os.path.exists(FIGURE_ID):
    os.makedirs(FIGURE_ID)

if not os.path.exists(DATA_ID):
    os.makedirs(DATA_ID)

def image_path(fig_id):
    return os.path.join(FIGURE_ID, fig_id)

def data_path(dat_id):
    return os.path.join(DATA_ID, dat_id)

def save_fig(fig_id):
    plt.savefig(image_path(fig_id) + ".png", format='png')

def R2(y_data, y_model):
    return 1 - np.sum((y_data - y_model) ** 2) / np.sum((y_data - np.mean(y_data)) ** 2)
def MSE(y_data,y_model):
    n = np.size(y_model)
    return np.sum((y_data-y_model)**2)/n

infile = open(data_path("EoS.csv"),'r')

# Read the EoS data as  csv file and organized into two arrays with density and energies
EoS = pd.read_csv(infile, names=('Density', 'Energy'))
EoS['Energy'] = pd.to_numeric(EoS['Energy'], errors='coerce')
EoS = EoS.dropna()
Energies = EoS['Energy']
Density = EoS['Density']
#  The design matrix now as function of various polytrops
X = np.zeros((len(Density),5))
X[:,0] = 1
X[:,1] = Density**(2.0/3.0)
X[:,2] = Density
X[:,3] = Density**(4.0/3.0)
X[:,4] = Density**(5.0/3.0)
# We split the data in test and training data
X_train, X_test, y_train, y_test = train_test_split(X, Energies, test_size=0.2)
# matrix inversion to find beta
beta = np.linalg.inv(X_train.T.dot(X_train)).dot(X_train.T).dot(y_train)
# and then make the prediction
ytilde = X_train @ beta
print("Training R2")
print(R2(y_train,ytilde))
print("Training MSE")
print(MSE(y_train,ytilde))
ypredict = X_test @ beta
print("Test R2")
print(R2(y_test,ypredict))
print("Test MSE")
print(MSE(y_test,ypredict))

\end{minted}
\end{frame}

\begin{frame}[plain,fragile]
\frametitle{Exercises}

Here are three possible exercises for weeks 34 and  35.

% --- begin exercise ---
\begin{doconceexercise}
\refstepcounter{doconceexercisecounter}

\exercisesection{Exercise \thedoconceexercisecounter: Setting up various Python environments}
                             

The first exercise here is of a mere technical art. We want you to have 
\begin{itemize}
\item git as a version control software and to establish a user account on a provider like GitHub. Other providers like GitLab etc are equally fine. You can also use the University of Oslo \href{{https://www.uio.no/tjenester/it/maskin/filer/versjonskontroll/github.html}}{GitHub facilities}. 

\item Install various Python packages
\end{itemize}

\noindent
We will make extensive use of Python as programming language and its
myriad of available libraries.  You will find
IPython/Jupyter notebooks invaluable in your work.  You can run \textbf{R}
codes in the Jupyter/IPython notebooks, with the immediate benefit of
visualizing your data. You can also use compiled languages like C++,
Rust, Fortran etc if you prefer. The focus in these lectures will be
on Python.

If you have Python installed (we recommend Python3) and you feel
pretty familiar with installing different packages, we recommend that
you install the following Python packages via \textbf{pip} as 

\begin{enumerate}
\item pip install numpy scipy matplotlib ipython scikit-learn sympy pandas pillow 
\end{enumerate}

\noindent
For \textbf{Tensorflow}, we recommend following the instructions in the text of 
\href{{http://shop.oreilly.com/product/0636920052289.do}}{Aurelien Geron, Hands‑On Machine Learning with Scikit‑Learn and TensorFlow, O'Reilly}

We will come back to \textbf{tensorflow} later. 

For Python3, replace \textbf{pip} with \textbf{pip3}.

For OSX users we recommend, after having installed Xcode, to
install \textbf{brew}. Brew allows for a seamless installation of additional
software via for example 

\begin{enumerate}
\item brew install python3
\end{enumerate}

\noindent
For Linux users, with its variety of distributions like for example the widely popular Ubuntu distribution,
you can use \textbf{pip} as well and simply install Python as 

\begin{enumerate}
\item sudo apt-get install python3  (or python for Python2.7)
\end{enumerate}

\noindent
If you don't want to perform these operations separately and venture
into the hassle of exploring how to set up dependencies and paths, we
recommend two widely used distrubutions which set up all relevant
dependencies for Python, namely 

\begin{itemize}
\item \href{{https://docs.anaconda.com/}}{Anaconda}, 
\end{itemize}

\noindent
which is an open source
distribution of the Python and R programming languages for large-scale
data processing, predictive analytics, and scientific computing, that
aims to simplify package management and deployment. Package versions
are managed by the package management system \textbf{conda}. 

\begin{itemize}
\item \href{{https://www.enthought.com/product/canopy/}}{Enthought canopy} 
\end{itemize}

\noindent
is a Python
distribution for scientific and analytic computing distribution and
analysis environment, available for free and under a commercial
license.

We recommend using \textbf{Anaconda} if you are not too familiar with setting paths in a terminal environment.

\end{doconceexercise}
% --- end exercise ---

% --- begin exercise ---
\begin{doconceexercise}
\refstepcounter{doconceexercisecounter}

\exercisesection{Exercise \thedoconceexercisecounter: making your own data and exploring scikit-learn}
                             

We will generate our own dataset for a function $y(x)$ where $x \in [0,1]$ and defined by random numbers computed with the uniform distribution. The function $y$ is a quadratic polynomial in $x$ with added stochastic noise according to the normal distribution $\cal {N}(0,1)$.
The following simple Python instructions define our $x$ and $y$ values (with 100 data points).



\begin{minted}[fontsize=\fontsize{9pt}{9pt},linenos=false,mathescape,baselinestretch=1.0,fontfamily=tt,xleftmargin=2mm]{python}
x = np.random.rand(100,1)
y = 2.0+5*x*x+0.1*np.random.randn(100,1)

\end{minted}


\begin{enumerate}
\item Write your own code (following the examples under the \href{{https://compphysics.github.io/MachineLearning/doc/LectureNotes/_build/html/chapter1.html}}{regression notes}) for computing the parametrization of the data set fitting a second-order polynomial. 

\item Use thereafter \textbf{scikit-learn} (see again the examples in the regression slides) and compare with your own code.   

\item Using scikit-learn, compute also the mean square error, a risk metric corresponding to the expected value of the squared (quadratic) error defined as
\end{enumerate}

\noindent
\[ MSE(\bm{y},\bm{\tilde{y}}) = \frac{1}{n}
\sum_{i=0}^{n-1}(y_i-\tilde{y}_i)^2, 
\] 
and the $R^2$ score function.
If $\tilde{\bm{y}}_i$ is the predicted value of the $i-th$ sample and $y_i$ is the corresponding true value, then the score $R^2$ is defined as
\[
R^2(\bm{y}, \tilde{\bm{y}}) = 1 - \frac{\sum_{i=0}^{n - 1} (y_i - \tilde{y}_i)^2}{\sum_{i=0}^{n - 1} (y_i - \bar{y})^2},
\]
where we have defined the mean value  of $\bm{y}$ as
\[
\bar{y} =  \frac{1}{n} \sum_{i=0}^{n - 1} y_i.
\]
You can use the functionality included in scikit-learn. If you feel for it, you can use your own program and define functions which compute the above two functions. 
Discuss the meaning of these results. Try also to vary the coefficient in front of the added stochastic noise term and discuss the quality of the fits.

% --- begin solution of exercise ---
\paragraph{Solution.}
The code here is an example of where we define our own design matrix and fit parameters $\beta$.









































\begin{minted}[fontsize=\fontsize{9pt}{9pt},linenos=false,mathescape,baselinestretch=1.0,fontfamily=tt,xleftmargin=2mm]{python}
import os
import numpy as np
import pandas as pd
import matplotlib.pyplot as plt
from sklearn.model_selection import train_test_split

def save_fig(fig_id):
    plt.savefig(image_path(fig_id) + ".png", format='png')

def R2(y_data, y_model):
    return 1 - np.sum((y_data - y_model) ** 2) / np.sum((y_data - np.mean(y_data)) ** 2)
def MSE(y_data,y_model):
    n = np.size(y_model)
    return np.sum((y_data-y_model)**2)/n

x = np.random.rand(100)
y = 2.0+5*x*x+0.1*np.random.randn(100)


#  The design matrix now as function of a given polynomial
X = np.zeros((len(x),3))
X[:,0] = 1.0
X[:,1] = x
X[:,2] = x**2
# We split the data in test and training data
X_train, X_test, y_train, y_test = train_test_split(X, y, test_size=0.2)
# matrix inversion to find beta
beta = np.linalg.inv(X_train.T @ X_train) @ X_train.T @ y_train
print(beta)
# and then make the prediction
ytilde = X_train @ beta
print("Training R2")
print(R2(y_train,ytilde))
print("Training MSE")
print(MSE(y_train,ytilde))
ypredict = X_test @ beta
print("Test R2")
print(R2(y_test,ypredict))
print("Test MSE")
print(MSE(y_test,ypredict))

\end{minted}


% --- end solution of exercise ---

\end{doconceexercise}
% --- end exercise ---

% --- begin exercise ---
\begin{doconceexercise}
\refstepcounter{doconceexercisecounter}

\exercisesection{Exercise \thedoconceexercisecounter: Normalizing our data}
                             

A much used approach before starting to train the data is  to preprocess our
data. Normally the data may need a rescaling and/or may be sensitive
to extreme values. Scaling the data renders our inputs much more
suitable for the algorithms we want to employ.

\textbf{Scikit-Learn} has several functions which allow us to rescale the
data, normally resulting in much better results in terms of various
accuracy scores.  The \textbf{StandardScaler} function in \textbf{Scikit-Learn}
ensures that for each feature/predictor we study the mean value is
zero and the variance is one (every column in the design/feature
matrix).  This scaling has the drawback that it does not ensure that
we have a particular maximum or minimum in our data set. Another
function included in \textbf{Scikit-Learn} is the \textbf{MinMaxScaler} which
ensures that all features are exactly between $0$ and $1$. The

The \textbf{Normalizer} scales each data
point such that the feature vector has a euclidean length of one. In other words, it
projects a data point on the circle (or sphere in the case of higher dimensions) with a
radius of 1. This means every data point is scaled by a different number (by the
inverse of it’s length).
This normalization is often used when only the direction (or angle) of the data matters,
not the length of the feature vector.

The \textbf{RobustScaler} works similarly to the StandardScaler in that it
ensures statistical properties for each feature that guarantee that
they are on the same scale. However, the RobustScaler uses the median
and quartiles, instead of mean and variance. This makes the
RobustScaler ignore data points that are very different from the rest
(like measurement errors). These odd data points are also called
outliers, and might often lead to trouble for other scaling
techniques.

It also common to split the data in a \textbf{training} set and a \textbf{testing} set. A typical split is to use $80\%$ of the data for training and the rest
for testing. This can be done as follows with our design matrix $\bm{X}$ and data $\bm{y}$ (remember to import \textbf{scikit-learn})



\begin{minted}[fontsize=\fontsize{9pt}{9pt},linenos=false,mathescape,baselinestretch=1.0,fontfamily=tt,xleftmargin=2mm]{python}
# split in training and test data
X_train, X_test, y_train, y_test = train_test_split(X,y,test_size=0.2)

\end{minted}

Then we can use the standard scaler to scale our data as





\begin{minted}[fontsize=\fontsize{9pt}{9pt},linenos=false,mathescape,baselinestretch=1.0,fontfamily=tt,xleftmargin=2mm]{python}
scaler = StandardScaler()
scaler.fit(X_train)
X_train_scaled = scaler.transform(X_train)
X_test_scaled = scaler.transform(X_test)

\end{minted}


In this exercise we want you to to compute the MSE for the training
data and the test data as function of the complexity of a polynomial,
that is the degree of a given polynomial. We want you also to compute the $R2$ score as function of the complexity of the model for both training data and test data.  You should also run the calculation with and without scaling. 

One of 
the aims is to reproduce Figure 2.11 of \href{{https://github.com/CompPhysics/MLErasmus/blob/master/doc/Textbooks/elementsstat.pdf}}{Hastie et al}.

Our data is defined by $x\in [-3,3]$ with a total of for example $100$ data points.







\begin{minted}[fontsize=\fontsize{9pt}{9pt},linenos=false,mathescape,baselinestretch=1.0,fontfamily=tt,xleftmargin=2mm]{python}
np.random.seed()
n = 100
maxdegree = 14
# Make data set.
x = np.linspace(-3, 3, n).reshape(-1, 1)
y = np.exp(-x**2) + 1.5 * np.exp(-(x-2)**2)+ np.random.normal(0, 0.1, x.shape)

\end{minted}

where $y$ is the function we want to fit with a given polynomial.

% --- begin subexercise ---
\subex{a)}
Write a first code which sets up a design matrix $X$ defined by a fifth-order polynomial.  Scale your data and split it in training and test data.

% --- end subexercise ---

% --- begin subexercise ---
\subex{b)}
Perform an ordinary least squares and compute the means squared error and the $R2$ factor for the training data and the test data, with and without scaling.

% --- end subexercise ---

% --- begin subexercise ---
\subex{c)}
Add now a model which allows you to make polynomials up to degree $15$.  Perform a standard OLS fitting of the training data and compute the MSE and $R2$ for the training and test data and plot both test and training data MSE and $R2$ as functions of the polynomial degree. Compare what you see with Figure 2.11 of Hastie et al.~Comment your results. For which polynomial degree do you find an optimal MSE (smallest value)?

% --- end subexercise ---

\end{doconceexercise}
% --- end exercise ---
\end{frame}

\end{document}
